% !TeX root = ../../Book3.tex
% 纲要标题：## 第27章*　Ricci 流、规范消除与奇点模型
% 纲要位置：几何学计划纲要.md，第 3240 行。
% 正文状态：初稿；尚未正式审核。

% 纲要内容开始
% 
% > 本章目标与证明范围：紧光滑初始度量的短时存在、唯一性和基本演化为完整模块。熵、非塌缩和手术按实际选取的深层结果区分证明与引用，不以二维算例替代高维理论。
% 
% 纲要内容结束

\OptionalChapter{Ricci 流、规范消除与奇点模型}{g3:ch:27}
\BookChapterNumber{g3:ch:27}

\begin{chapterabstract}
Ricci 流以 Ricci 张量改变度量，使曲率满足带二次反应项的扩散方程。本章从微分同胚不变性导致的主符号退化出发，用 DeTurck 向量场构造严格抛物系统，证明闭光滑初始度量的短时存在、唯一性及曲率有界时的延拓。随后推导体积、标量曲率、Ricci 张量和完整曲率算子的演化，证明任意维非负曲率算子的保持与正时间曲率导数估计。奇点部分讨论抛物重缩放、梯度孤立子以及 Perelman 两个泛函的单调公式，精确引用有限时间非塌缩定理。曲面、体积归一化和 Nil 齐性模型提供具体计算，三维手术与几何化则在末节说明所用的外部结果。
\end{chapterabstract}

\begin{learninggoals}
\begin{itemize}
\item 计算 Ricci 流的退化主符号，并说明 DeTurck 项怎样恢复强抛物性。
\item 从第26章的拟线性理论构造流，通过时间依赖微分同胚证明存在与唯一性。
\item 推导完整曲率演化，区分标量比较、导数估计与曲率锥的保持。
\item 解释有限时曲率界为何足以恢复光滑终点度量并继续求解。
\item 计算孤立子与抛物放大，证明 \(\mathcal F\)、\(\mathcal W\) 单调公式，准确使用非塌缩条件。
\item 求解曲面和齐性模型，区分光滑延拓、体积归一化与三维手术。
\end{itemize}
\end{learninggoals}

单位圆球的 Ricci 张量为 \((n-1)g\)。若令度量按 \(\partial_tg=-2\operatorname{Ric}\) 演化，圆球仍然是圆球，其半径平方以常速度下降；平坦度量则保持不变。对一般初值，曲率在空间中扩散，却又受到二次反应项的影响，因此既可能变得均匀，也可能在有限时间内集中。Hamilton 引入的 Ricci 流把这些变化组织成一个几何初值问题。\footnote{Richard Streit Hamilton（1943--2024），美国数学家，研究微分几何与几何分析，创立 Ricci 流方法；姓名及生卒年见\href{https://www.math.columbia.edu/2024/10/02/richard-streit-hamilton-1943-2024/}{Columbia University 数学系纪念文章}。}

\ProseParagraphBreak
要开始求解，首先必须处理方程的退化。度量可以沿微分同胚拉回，表示同一个几何；这种自由在最高阶符号中留下了核。第25章已经在曲率线性化中看到同一现象。本章加入由初始度量确定的 DeTurck 项，再用第26章的强抛物理论求解，最后沿一个时间依赖向量场拉回。唯一性需要逆向比较：从任意 Ricci 流构造相同的严格抛物规范，才能证明原来的两个参数化解相等。

\ProseParagraphBreak
局部求解以后，困难转向先验估计。标量曲率满足一个容易比较的方程，完整曲率却包含多个方向的耦合。固定一个正交向量丛来追踪曲率算子，可以把“非负”写成固定凸锥中的条件。曲率的协变导数估计进一步说明，有限时间内只要完整曲率保持有界，初始短时构造就可以重新启动。相反，奇点处必须放大曲率尺度；孤立子、单调泛函和非塌缩为识别放大后的几何提供不同的信息。

\ProseParagraphBreak
本章使用\secref{g3:sec:2501}的曲率变分、\secref{g3:sec:2601}的最大值原理及\secref{g3:sec:2602}的拟线性短时理论。曲率张量与曲率算子的约定沿用第3章，曲面面积公式使用第17章的二维 Gauss--Bonnet 公式。短时理论、曲率保持及两个单调公式在正文证明；非塌缩和三维手术存在性在相应陈述处明确引用。

% !TeX root = ../../Book3.tex
% 纲要标题：### 27.1 Ricci 流与 DeTurck 规范
% 纲要位置：几何学计划纲要.md，第 3244 行。
% 正文状态：初稿；尚未正式审核。

% 纲要内容开始
% 
% * 方程、缩放和微分同胚不变性
% * 线性化中的非抛物方向
% * DeTurck 向量场与严格抛物方程
% 
% 纲要内容结束

\section{Ricci 流与 DeTurck 规范}
\label{g3:sec:2701}

沿曲率调整度量，希望同时改变各个方向的长度。Ricci 张量与度量都是对称二张量，因此可以直接把前者作为后者的时间导数。不过，坐标变换也会改变度量分量，曲率方程的最高阶项必须容纳这种自由。本节先明确方程的缩放，再利用第25章的线性化找出退化，最后加入一个由背景联络确定的 Lie 导数，使方程能够使用第26章的强抛物理论。

\subsection{方程、缩放和微分同胚不变性}
\label{g3:subsec:2701:01}
\begin{definition}\label{g3:def:ricci-flow}
	设 \(M^n\) 为光滑无边界流形，\(n\ge2\)，\(I\) 为时间区间。联合光滑的黎曼度量族 \(g(t)\)、\(t\in I\)，若满足
	\begin{equation}\label{g3:eq:ricci-flow}
		\partial_tg=-2\operatorname{Ric}_{g(t)},
	\end{equation}
	则称为 \(M\) 上的\term[Ricci liu]{Ricci 流}[Ricci flow]。求导时固定流形上的点和切向量，右侧使用当前度量的 Ricci 协变二张量。
\end{definition}

在正交特征方向 \(v\) 上，\(\partial_tg(v,v)=-2\operatorname{Ric}(v,v)\)：正 Ricci 方向缩短，负 Ricci 方向增长。这里调整的是内在度量，并未把 \(M\) 放入一个环境空间。常数 \(-2\) 使消除坐标退化后的最高阶项具有通常热方程的系数。

\begin{proposition}\label{g3:prop:ricci-flow-invariance}
	设 \(M^n\) 为光滑无边界流形，\(n\ge2\)，\(g(t)\) 是 \(M\) 上的光滑 Ricci 流，\(c>0\)，\(\psi:M\to M\) 为时间不变的微分同胚。则 \(\psi^*g(t)\) 与 \(g_c(t)=c\,g(t/c)\) 在各自定义的时间区间上也满足 Ricci 流。对每个光滑度量 \(g\)，有
	\begin{equation}\label{g3:eq:ricci-flow-scaling}
		\operatorname{Ric}_{cg}=\operatorname{Ric}_g,\quad
		s_{cg}=c^{-1}s_g,\quad
		|\mathrm{Rm}_{cg}|_{cg}=c^{-1}|\mathrm{Rm}_g|_g,\quad
		\mu_{cg}=c^{n/2}\mu_g.
	\end{equation}
\end{proposition}
\begin{proof}
	Levi--Civita 联络由度量自然确定，所以拉回联络的曲率收缩给 \(\operatorname{Ric}_{\psi^*g}=\psi^*\operatorname{Ric}_g\)。时间求导与固定拉回交换，得到第一种不变性。
	常数缩放使 Koszul 公式的两边同乘 \(c\)，故联络以及 \((1,3)\) 型曲率不变。Ricci 是这种曲率的自然迹，因而也不变。标量曲率再用逆度量取迹，所以乘 \(c^{-1}\)。协变四阶曲率张量乘 \(c\)，其平方范数中的四个逆度量共乘 \(c^{-4}\)，故范数乘 \(c^{-1}\)。体积密度由行列式平方根给出最后一式。最后
	\(\partial_t(cg(t/c))=-2\operatorname{Ric}_{g(t/c)}=-2\operatorname{Ric}_{g_c(t)}\)，证明时间与度量的联合缩放。
\end{proof}

\begin{example}\label{g3:exm:ricci-flow-einstein-scaling}
	若 \(\operatorname{Ric}_{g_0}=\lambda g_0\)，其中 \(\lambda\) 为常数，则
	\[
		g(t)=(1-2\lambda t)g_0
	\]
	在 \(1-2\lambda t>0\) 的区间上求解。直接验证时必须保留 \(\operatorname{Ric}_{g(t)}=\lambda g_0\)，而相对于当前度量的 Einstein 常数是 \(\lambda/(1-2\lambda t)\)。正、零、负 \(\lambda\) 分别给出收缩、静止与膨胀。特别地，单位圆球的解为 \((1-2(n-1)t)g_{S^n}\)，其截面曲率为 \((1-2(n-1)t)^{-1}\)。
\end{example}

时间依赖的微分同胚会产生额外项。若 \(\phi_t\) 是光滑微分同胚族，\(X_t=\partial_t\phi_t\circ\phi_t^{-1}\)，则
\begin{equation}\label{g3:eq:ricci-flow-time-pullback}
	\partial_t(\phi_t^*g(t))
	=\phi_t^*\bigl(\partial_tg+\mathcal L_{X_t}g\bigr).
\end{equation}
确实，在坐标式 \((\phi_t^*g)_{ij}=g_{ab}(\phi_t,t)\partial_i\phi_t^a\partial_j\phi_t^b\) 中求导，移动取值点与两个 Jacobian 因子的导数恰给 Lie 导数的三项。这一公式将把加入的规范项与参数运动相互转换。

\subsection{线性化中的非抛物方向}
\label{g3:subsec:2701:02}
固定一个光滑度量 \(g\)，把 \(h\) 看作度量的变分。\Cref{g3:thm:ricci-scalar-linearization}已经在任意维局部计算中证明
\[
	\mathrm D(-2\operatorname{Ric})_g(h)
	=-\Delta_Lh+2\delta_g^*\beta_gh,
	\qquad \beta_gh=\delta_gh+\tfrac12\mathrm d\operatorname{tr}_gh.
\]
这里 \((\delta h)_j=-\nabla^ih_{ij}\)，\(2\delta^*\omega=\nabla_i\omega_j+\nabla_j\omega_i\)，而 \(\Delta_L\) 的主部是非负粗 Laplace 算子。该公式的推导不使用第25章讨论 Einstein 模空间时的 \(n\ge3\) 限制，所以也适用于曲面。

\begin{proposition}\label{g3:prop:ricci-flow-symbol-degeneracy}
	设 \((M^n,g)\) 为光滑无边界黎曼流形，\(n\ge2\)，\(h\) 为光滑对称协变二张量。对非零实余向量 \(\xi\)，按 \(\partial_j\mapsto i\xi_j\) 取主符号，则
	\begin{equation}\label{g3:eq:ricci-flow-degenerate-symbol}
		\bigl[\sigma_2(\mathrm D(-2\operatorname{Ric})_g)(\xi)h\bigr]_{ij}
		=-|\xi|^2h_{ij}+\xi_i\xi^ah_{aj}+\xi_j\xi^ah_{ai}
		-\xi_i\xi_j\operatorname{tr}_gh.
	\end{equation}
	对任意余向量 \(\eta\)，张量 \(h=\xi\otimes\eta+\eta\otimes\xi\) 都属于此符号的核。因此 Ricci 参数方程不是对全部度量分量的强抛物系统。
\end{proposition}
\begin{proof}
	在一点取正交坐标。\(-\Delta_L\) 的符号为 \(-|\xi|^2h\)，而 \(\beta h\) 的符号为 \(-i\xi^ah_{aj}+\tfrac i2\xi_j\operatorname{tr}h\)。再作用 \(2\delta^*\) 就得到所示三项。
	对给定 \(h\)，有 \(\xi^ah_{aj}=|\xi|^2\eta_j+\xi_j\langle\xi,\eta\rangle\) 以及 \(\operatorname{tr}h=2\langle\xi,\eta\rangle\)。代入后全部抵消。非零 \(\eta\) 给出非零 \(h\)，所以不存在强抛物性要求的严格符号下界。
\end{proof}

核中的张量正是 \(\mathcal L_Xg\) 的一阶符号形式。这种退化也可从自然性看出：若 \(g_0\) 为 Ricci 平坦度量，任意固定微分同胚拉回的度量仍 Ricci 平坦；直接以度量分量衡量变化，就包含大量只改变点的标记的方向。规范的任务是给这些方向补上确定的演化。

\subsection{DeTurck 向量场与严格抛物方程}
\label{g3:subsec:2701:03}
固定一个光滑背景度量 \(\bar g\)，令 \(\bar\nabla\) 为其 Levi--Civita 联络。两条联络之差
\[
	C^k_{ij}=\Gamma(g)^k_{ij}-\Gamma(\bar g)^k_{ij}
	=\frac12g^{ka}\bigl(\bar\nabla_i g_{ja}+\bar\nabla_j g_{ia}-\bar\nabla_a g_{ij}\bigr)
\]
是张量。于是可以定义全局向量场
\begin{equation}\label{g3:eq:deturck-vector-field}
	W(g,\bar g)^k\coloneqq g^{ij}C^k_{ij}.
\end{equation}
\symidx{\(W(g,\bar g)\)：相对于背景度量 \(\bar g\) 的 DeTurck 向量场}
它称为\term[DeTurck xiangliangchang]{DeTurck 向量场}[DeTurck vector field]。\footnote{Dennis M. DeTurck（1954年生），美国数学家，研究偏微分方程、黎曼几何及给定 Ricci 曲率问题；其规范方法用于把 Ricci 流化为强抛物方程。}相应的\term[Ricci DeTurck liu]{Ricci--DeTurck 流}[Ricci--DeTurck flow]为
\begin{equation}\label{g3:eq:ricci-deturck-flow}
	\partial_tg=-2\operatorname{Ric}_g+\mathcal L_{W(g,\bar g)}g.
\end{equation}

\begin{proposition}\label{g3:prop:ricci-deturck-parabolic}
	设 \(M\) 为光滑无边界流形，\(\bar g\) 为固定光滑黎曼度量。记 DeTurck 场 \(W(g,\bar g)^k=g^{ij}(\Gamma(g)^k_{ij}-\Gamma(\bar g)^k_{ij})\)，其中 \(\Gamma\) 为相应 Levi--Civita 联络系数。所讨论方程为 \(\partial_tg=-2\operatorname{Ric}_g+\mathcal L_{W(g,\bar g)}g\)，\(\bar\nabla\) 为 \(\bar g\) 的联络。
	在正定对称二张量的开集上，式\eqref{g3:eq:ricci-deturck-flow}可写为
	\begin{equation}\label{g3:eq:ricci-deturck-principal}
		\partial_tg_{ij}
		=g^{ab}\bar\nabla_a\bar\nabla_b g_{ij}
		+\mathcal Q_{ij}(x,g,\bar\nabla g),
	\end{equation}
	其中 \(\mathcal Q\) 对所示变量光滑。其主部在 \(\operatorname{Sym}^2\mathrm T^*M\) 上为标量系数乘恒等算子；在 \(c^{-1}\bar g\le g\le c\bar g\) 的区域上一致强抛物。
\end{proposition}
\begin{proof}
	先只保留局部坐标中的二阶度量导数，记所有至多一阶的项为 \(Q\)。从联络曲率公式得到
	\[
		2\operatorname{Ric}_{ij}
		=g^{ab}\bigl(\partial_a\partial_i g_{bj}
		+\partial_a\partial_j g_{bi}-\partial_a\partial_b g_{ij}
		-\partial_i\partial_j g_{ab}\bigr)+Q_{ij}.
	\]
	降低 \(W\) 的指标后，含度量一阶导数的部分为
	\(W_j=g^{ab}(\partial_a g_{bj}-\tfrac12\partial_jg_{ab})\)；背景联络项只依赖 \(x,g\)。因此 Lie 导数的二阶部分为
	\[
		(\mathcal L_Wg)_{ij}
		=g^{ab}\bigl(\partial_i\partial_a g_{bj}
		+\partial_j\partial_a g_{bi}-\partial_i\partial_jg_{ab}\bigr)+Q'_{ij}.
	\]
	将两式按 \(-2\operatorname{Ric}+\mathcal L_Wg\) 相加，混合二阶导数相消，只留下 \(g^{ab}\partial_a\partial_b g_{ij}\)。将普通导数换成固定背景的协变导数，只改变光滑低阶项，得到式\eqref{g3:eq:ricci-deturck-principal}。
	若 \(\xi\) 为余向量、\(v\) 为对称二张量，主二次型是 \(g^{ab}\xi_a\xi_b|v|_{\bar g}^2\)，在所给区域至少为 \(c^{-1}|\xi|_{\bar g}^2|v|_{\bar g}^2\)。这符合\cref{g3:def:quasilinear-strong-parabolicity}的定义。
\end{proof}

还可在线性层面核对所选符号。在 \(g=\bar g\) 处，\(W=0\)，联络变分公式给
\[
	\bigl(\mathrm DW_{\bar g}(h)\bigr)^\flat
	=\nabla^ih_{ij}-\tfrac12\nabla_j\operatorname{tr}h
	=-\beta_{\bar g}h.
\]
所以 \(\mathrm D(\mathcal L_Wg)=-2\delta^*\beta h\)，恰好消去 Ricci 线性化中的退化项，剩下 \(-\Delta_Lh\)。这个计算也固定了下一节拉回微分同胚必须采用 \(-W\) 作为速度的符号。
% 纲要条目排为编号小节标题，并进入目录与书签；说明段仍保留于纲要注释中。

% !TeX root = ../../Book3.tex
% 纲要标题：### 27.2 短时理论的证明
% 纲要位置：几何学计划纲要.md，第 3250 行。
% 正文状态：初稿；尚未正式审核。

% 纲要内容开始
% 
% * 调用第26章的拟线性存在
% * 时间依赖微分同胚的构造
% * 拉回恢复 Ricci 流
% * 唯一性的规范比较
% * 曲率界下的延拓判据
% 
% 纲要内容结束

\section{短时理论的证明}
\label{g3:sec:2702}

上一节已经把规范固定后的最高阶部分化为正定逆度量乘以二阶空间导数。现在需要把这个符号计算落实成一个光滑度量族，再去掉规范项。存在性与唯一性沿相反方向使用这一步：构造时先求规范方程，再作拉回；比较两个已知 Ricci 流时，则先用调和映射热流把它们放进同一个规范。最后从曲率控制恢复固定坐标中的度量控制，说明什么会阻止光滑解继续存在。

\subsection{调用第26章的拟线性存在}
\label{g3:subsec:2702:01}
固定闭光滑流形 \(M^n\)、\(n\ge2\)，光滑初始度量 \(g_0\) 及光滑背景度量 \(\bar g\)。本节用 \(h\) 表示规范方程的未知度量，把 \(g\) 留给恢复出的 Ricci 流。背景联络 \(\bar\nabla\) 及全部 Banach 空间都固定不动。

\begin{proposition}\label{g3:prop:ricci-deturck-short-time}
 设 \(M^n\) 闭，\(n\ge2\)，\(g_0,\bar g\) 为光滑黎曼度量。则存在 \(T>0\) 及唯一的光滑正定度量族 \(h(t)\)，满足
 \[
  \partial_th=-2\operatorname{Ric}_h+\mathcal L_{W(h,\bar g)}h,
  \qquad h(0)=g_0.
 \]
 唯一性在光滑正定解的共同存在区间内成立，光滑性包括初时刻的单侧导数。
\end{proposition}
\begin{proof}
 取固定向量丛 \(E=\operatorname{Sym}^2\mathrm T^*M\)，配上 \(\bar g\) 诱导的丛度量与联络。\Cref{g3:prop:ricci-deturck-parabolic}把方程写成
 \[
  \partial_th=h^{ab}\bar\nabla_a\bar\nabla_bh
       +Q(x,h,\bar\nabla h),
 \]
 其中 \(Q\) 在 \(h\) 正定的一阶喷射开集上光滑。对 \(\xi\in\mathrm T_x^*M\)、\(v\in E_x\)，主符号的二次型为
 \[
  h^{ab}\xi_a\xi_b|v|_{\bar g}^2.
 \]
 在初始喷射图上，紧性和 \(g_0\) 正定给出它相对于 \(|\xi|_{\bar g}^2|v|_{\bar g}^2\) 的统一正下界；缩小喷射邻域仍保留正下界。取整数 \(s>n/2+3\)，\cref{g3:thm:quasilinear-parabolic-short-time}的全部假设遂满足。该定理给出光滑短时解；\(C^0\) 连续性保证缩短时间后仍正定。其唯一性由\cref{g3:prop:quasilinear-parabolic-uniqueness}给出。
\end{proof}

这里无需初始度量接近背景度量。两者在紧流形上总可作统一比较，强抛物常数可以依赖这组给定数据。正定性是未知量所处的开放条件，不能只看各坐标分量有界而忽略它。

\subsection{时间依赖微分同胚的构造}
\label{g3:subsec:2702:02}
一旦得到光滑的 \(h(t)\)，向量场 \(W(h(t),\bar g)\) 就成为已知的光滑时变系数。去规范所需的映射由常微分方程构造，不再是第二个偏微分方程问题。

\begin{lemma}\label{g3:lem:ricci-gauge-diffeomorphisms}
 设 \(M\) 闭，\(h(t)\) 是 \([0,T]\) 上的光滑正定度量族。则
 \begin{equation}\label{g3:eq:ricci-gauge-ode}
  \frac{\mathrm d}{\mathrm dt}\phi_t(x)
   =-W(h(t),\bar g)(\phi_t(x)),\qquad\phi_0(x)=x
 \end{equation}
 有唯一光滑解族 \(\phi_t:M\to M\)，每个 \(\phi_t\) 都是微分同胚。
\end{lemma}
\begin{proof}
 光滑时变向量场在有限坐标片内满足常微分方程的局部存在、唯一性和光滑依赖条件。闭性使解曲线在有限时间中不能逃出紧集；用有限坐标覆盖及紧时间区间上的系数界逐段延续，就得到整个 \([0,T]\) 上的解。给定终点 \(y\) 和时刻 \(t\)，反向求解同一常微分方程至时间零，得到 \(\phi_t\) 的逆。常微分方程的光滑依赖同时给出逆映射的光滑性。
\end{proof}

拉回一个本身随时间变化的张量时，既要微分张量，也要微分拉回映射。下面固定这个计算的方向与符号。

\begin{lemma}\label{g3:lem:time-dependent-pullback}
 设 \(\phi_t\) 是光滑微分同胚族，\(X_t=\dot\phi_t\circ\phi_t^{-1}\)，\(A(t)\) 是任意光滑张量族。则
 \begin{equation}\label{g3:eq:time-dependent-pullback}
  \frac{\mathrm d}{\mathrm dt}(\phi_t^*A(t))
  =\phi_t^*(\partial_tA+\mathcal L_{X_t}A).
 \end{equation}
\end{lemma}
\begin{proof}
 对函数，这就是链式法则。对二协变张量，在坐标中写
 \((\phi_t^*A)_{ij}=A_{ab}(\phi_t,t)\partial_i\phi_t^a\partial_j\phi_t^b\)。微分右端，把 \(\dot\phi_t=X_t\circ\phi_t\) 代入，括号中恰出现
 \[
  \partial_tA_{ab}+X_t^c\partial_cA_{ab}
   +A_{cb}\partial_aX_t^c+A_{ac}\partial_bX_t^c,
 \]
 即 \(\partial_tA+\mathcal L_{X_t}A\) 的分量。对向量、一般协变和逆变张量按相同的链式法则计算，或使用张量积与缩并相容性，就得到一般公式。
\end{proof}

\subsection{拉回恢复 Ricci 流}
\label{g3:subsec:2702:03}
式\eqref{g3:eq:ricci-gauge-ode}中的负号恰好抵消规范方程中的正 Lie 导数。抵消之后，Ricci 张量在拉回下的自然性使方程仍用新的度量自身表达。

\begin{proposition}\label{g3:prop:ricci-flow-existence}
 设 \(M^n\) 闭，\(n\ge2\)，\(g_0\) 为光滑黎曼度量。取任意光滑背景度量 \(\bar g\)，令 \(h\) 为\cref{g3:prop:ricci-deturck-short-time}的解，\(\phi_t\) 为式\eqref{g3:eq:ricci-gauge-ode}的解。则
 \begin{equation}\label{g3:eq:ricci-flow-from-deturck}
  g(t)=\phi_t^*h(t)
 \end{equation}
 是以 \(g_0\) 为初值的光滑短时 Ricci 流。
\end{proposition}
\begin{proof}
 拉回保持正定性，\(\phi_0=\operatorname{id}\) 给出初值。将 \(X_t=-W(h(t),\bar g)\) 代入拉回公式，得到
 \[
  \partial_tg
  =\phi_t^*\bigl(-2\operatorname{Ric}_h
       +\mathcal L_Wh-\mathcal L_Wh\bigr)
  =-2\operatorname{Ric}_{\phi_t^*h}.
 \]
 最后一个等号来自 Levi--Civita 联络、曲率及缩并在微分同胚下的相容性，因此就是 \(\partial_tg=-2\operatorname{Ric}_g\)。
\end{proof}

\begin{example}\label{g3:exm:ricci-time-dependent-flat-pullback}
 在平坦环面 \((\mathbb R/2\pi\mathbb Z)^n\)、\(n\ge2\) 上，取
 \[
  g_0=\sum_i(\mathrm dx_i)^2,\qquad
  \psi_t(x_1,x_2,\ldots)=(x_1+t\sin x_2,x_2,\ldots).
 \]
 每个 \(\psi_t\) 都是微分同胚，且
 \[
  \psi_t^*g_0=(\mathrm dx_1+t\cos x_2\,\mathrm dx_2)^2
       +\sum_{i=2}^n(\mathrm dx_i)^2
 \]
 始终平坦。但是它的时间导数不恒为零，故这条度量曲线不满足 Ricci 流方程。这个例子说明，时间依赖的坐标变换产生真实的 Lie 导数项；恢复 Ricci 流时，必须使用刚才所确定的向量场及符号。
\end{example}

\subsection{唯一性的规范比较}
\label{g3:subsec:2702:04}
要比较两个 Ricci 流，需要把每一个已知解都转成相同背景下的 Ricci--DeTurck 解。此时规范度量尚未知，不能先把它的 \(W\) 当作给定系数。调和映射热流则以给定的 Ricci 流作源度量，直接提供所需坐标。

\begin{lemma}\label{g3:lem:time-dependent-harmonic-gauge}
 设 \(M\) 闭，\(g(t)\)、\(0\le t\le T_0\)，为任意光滑正定度量族，\(\bar g\) 为固定光滑度量。则对充分小的 \(T>0\)，存在唯一光滑映射族
 \begin{equation}\label{g3:eq:ricci-harmonic-gauge-heat}
  F_t=\tau_{g(t),\bar g}(F),\qquad F(\cdot,0)=\operatorname{id}_M.
 \end{equation}
 每个 \(F(t)\) 是微分同胚。
\end{lemma}
\begin{proof}
 对源坐标 \(x^i\) 与目标坐标 \(y^\alpha\)，方程是
 \[
  \partial_tF^\alpha=g^{ij}(t)\bigl(
   \partial_i\partial_jF^\alpha
   -\Gamma^k_{ij}(g(t))\partial_kF^\alpha
   +\bar\Gamma^\alpha_{\beta\gamma}(F)
           \partial_iF^\beta\partial_jF^\gamma\bigr).
 \]
 用固定 \(\bar g\) 的指数映射写 \(F(x,t)=\exp_x^{\bar g}v(x,t)\)，其中 \(v\) 是固定丛 \(\mathrm TM\) 的小截面。时间微分与二阶空间微分的最高阶项都带同一个可逆因子 \(\mathrm D_v\exp_x^{\bar g}\)。消去该因子后，方程成为
 \[
  v_t=g^{ij}(t)\bar\nabla_i\bar\nabla_jv
       +B(t,x,v,\bar\nabla v),\qquad v(0)=0,
 \]
 其中 \(B\) 光滑。紧性和 \(g(0)\) 正定保证短时间内的一致强抛物性。

 \Cref{g3:thm:quasilinear-parabolic-short-time}的证明同样适用于这里的光滑时变系数，所需调整可直接核对。把算子写为 \(\mathscr P(v)=v_t-Q(t,v)\)，在时间常值初值处的线性化成为时变强抛物算子；\cref{g3:thm:linear-parabolic-isomorphism}已经允许这样的系数，并给出小时间区间上一致有界的逆。初始误差满足
 \[
  \|Q(t,0)\|_{L^2(0,T;H^{s-1})}
  \le\sqrt T\sup_{0\le t\le T_0}\|Q(t,0)\|_{H^{s-1}}.
 \]
 非线性映射的乘积、复合及微分估计对这组光滑时变系数仍一致成立。因此先选小球、再缩短时间的压缩构造不变。\Cref{g3:lem:parabolic-parameter-regularity}中的空间坐标变化保持时间变量，变换后的系数仍对时间和参数光滑；同一参数论证给出所有空间导数，方程再给出所有时间导数。两解之差的线性能量估计也不变，故得到光滑存在和唯一性。

 最后，\(F(t)\to\operatorname{id}\) 于 \(C^1\)。充分接近恒等映射的映射为微分同胚：有限坐标片中的逆函数估计保证邻近点的像不同，远离对角线的点对则由紧性与一致接近保证像不同；其像在每个分支中开且闭，故满射。缩短 \(T\) 就得到所需结论。
\end{proof}

\begin{lemma}\label{g3:lem:ricci-flow-to-deturck}
 设 \(g(t)\) 是闭流形上的光滑 Ricci 流，\(F(t)\) 为上一引理构造的微分同胚。定义目标上的度量及速度场
 \[
  h(t)=(F(t)^{-1})^*g(t),\qquad
  Y_t=\partial_tF\circ F^{-1}.
 \]
 则 \(Y_t=-W(h(t),\bar g)\)，且 \(h\) 满足以 \(g(0)\) 为初值的 Ricci--DeTurck 方程。
\end{lemma}
\begin{proof}
 由定义，\(F:(M,g)\to(M,h)\) 是等距映射。在 \(y=F(x)\) 所给的源坐标中，映射 \(F\) 的坐标表达为恒等映射，源度量的分量为 \(h\)，而热流的目标度量仍为 \(\bar g\)。因此张力场的坐标式直接给出
 \[
  Y_t^k=h^{ij}\bigl(\bar\Gamma^k_{ij}-\Gamma^k_{ij}(h)\bigr)
      =-W(h,\bar g)^k.
 \]
 另一方面 \(g=F^*h\)，应用式\eqref{g3:eq:time-dependent-pullback}，并对两边作 \(F_*=(F^{-1})^*\)，得到
 \[
  -2\operatorname{Ric}_h=\partial_th+\mathcal L_{Y_t}h.
 \]
 代入 \(Y_t=-W(h,\bar g)\) 即得规范方程。初值由 \(F(0)=\operatorname{id}\) 确定。
\end{proof}

\begin{theorem}\label{g3:thm:ricci-flow-short-time}
 设 \(M^n\) 为闭光滑流形，\(n\ge2\)，\(g_0\) 为光滑黎曼度量。则存在以 \(g_0\) 为初值的光滑短时 Ricci 流；任意两个具有该初值的光滑 Ricci 流在共同存在区间内相同。因此存在唯一的最大光滑解，定义在 \([0,T_{\max})\)，其中 \(0<T_{\max}\le\infty\)。
\end{theorem}
\begin{proof}
 存在性已经由\cref{g3:prop:ricci-flow-existence}证明。设 \(g_1,g_2\) 是同初值的两个解，选同一个 \(\bar g\)。分别求出式\eqref{g3:eq:ricci-harmonic-gauge-heat}的解 \(F_1,F_2\)，并令 \(h_i=(F_i^{-1})^*g_i\)。上一引理使两个 \(h_i\) 满足同一个 Ricci--DeTurck 初值问题，因此在短时间内 \(h_1=h_2=h\)。

 再由 \(\partial_tF_i\circ F_i^{-1}=-W(h,\bar g)\)，两映射都满足同一个常微分方程
 \[
  \partial_tF_i=-W(h,\bar g)\circ F_i,\qquad F_i(0)=\operatorname{id}.
 \]
 常微分方程唯一性给出 \(F_1=F_2\)，于是 \(g_1=F_1^*h=F_2^*h=g_2\)。只要两个解在某个时刻相同，就可在该时刻重新作同一短时比较。连续性与向右延续遂给出整个共同存在区间的唯一性。

 取所有光滑存在区间的并。共同区间上的唯一性使各解能够一致粘合；局部短时存在保证其时间长度的上确界为正。得到的解就是所述最大解。不同背景度量构造出的解也因这一唯一性而相同。
\end{proof}

\subsection{曲率界下的延拓判据}
\label{g3:subsec:2702:05}
短时存在定理使用固定坐标中的度量系数及其导数。几何上较自然的假设却是曲率有界；两者之间需要一段转换。本小节使用\cref{g3:thm:ricci-flow-curvature-derivatives}：闭 Ricci 流在有限时间区间上若 \(|\operatorname{Rm}|\le K\)，则对每个 \(a>0\) 和整数 \(r\ge1\)，\(\nabla^r\operatorname{Rm}\) 在 \(a\le t<T\) 上一致有界。该结果将在\secref{g3:sec:2703}由曲率演化和最大值原理证明，其证明不使用本小节的延拓结论。

\begin{lemma}\label{g3:lem:ricci-fixed-coordinate-control}
 设 \(g(t)\)、\(0\le t<T<\infty\)，为闭流形 \(M^n\) 上的光滑 Ricci 流，并且 \(\sup_{M\times[0,T)}|\operatorname{Rm}|_{g(t)}\le K\)。固定 \(0<a<T\)，令 \(g_a=g(a)\)、\(D=\nabla^{g_a}\)。则 \(g(t)\) 与 \(g_a\) 在 \([a,T)\) 上一致等价，而且对每个整数 \(r\ge0\)，
 \[
  \sup_{a\le t<T}\bigl(
   \|D^rg(t)\|_{C^0,g_a}+\|D^r\partial_tg(t)\|_{C^0,g_a}
  \bigr)<\infty.
 \]
\end{lemma}
\begin{proof}
 曲率界给出 \(|\operatorname{Ric}|_g\le c_nK\)。对固定非零切向量 \(v\)，流方程给出
 \[
  \left|\frac{\mathrm d}{\mathrm dt}\log g(t)(v,v)\right|
  \le2c_nK.
 \]
 从 \(a\) 积分到 \(t\)，得到
 \begin{equation}\label{g3:eq:ricci-metric-equivalence}
  e^{-2c_nK(T-a)}g_a\le g(t)\le e^{2c_nK(T-a)}g_a.
 \end{equation}
 所有固定阶数张量的 \(g(t)\) 范数与 \(g_a\) 范数因此可统一比较。

 令 \(S(t)=\nabla^{g(t)}-D\)，它是 \((1,2)\) 型张量，\(S(a)=0\)。由\cref{g3:lem:levi-civita-variation}对 \(\partial_tg=-2\operatorname{Ric}\) 的应用，
 \begin{equation}\label{g3:eq:ricci-connection-time-derivative}
  (\partial_tS)^k_{ij}
  =-g^{k\ell}\bigl(
    \nabla_i\operatorname{Ric}_{j\ell}
   +\nabla_j\operatorname{Ric}_{i\ell}
   -\nabla_\ell\operatorname{Ric}_{ij}\bigr).
 \end{equation}
 这里的 \(\nabla\) 使用正在演化的度量。曲率导数估计和范数比较使右端在 \([a,T)\) 上一致有界，时间积分遂给 \(S\) 的固定背景范数界。

 下面逐阶控制 \(D^rS\)，同时核对导数计数。对任意张量 \(A\)，有 \(DA=\nabla A+S*A\)，其中星号表示由张量类型确定的有限个缩并和正负号。反复使用此式，\(D^rA\) 可写成 \(\nabla^rA\) 加上若干乘积；每个乘积只含 \(\nabla^jA\)、\(j<r\)，以及 \(D^\ell S\)、\(0\le\ell\le r-1\)。这一描述由对等式再施一次 \(D\) 归纳得到：导数落到 \(S\) 因子上，阶数增加一；落到 \(\nabla^jA\) 上，则变成 \(\nabla^{j+1}A+S*\nabla^jA\)。

 取 \(A=\partial_tS\)。由式\eqref{g3:eq:ricci-connection-time-derivative}和 \(\nabla g=0\)，\(\nabla^jA\) 是 \(g^{-1}\) 与 \(\nabla^{j+1}\operatorname{Ric}\) 的固定缩并，均已由曲率导数估计控制。因此，若 \(D^\ell S\)、\(\ell<r\)，已经有界，就得到 \(D^r\partial_tS\) 的一致界。由于 \(D\) 不随时间变化，
 \[
  D^rS(t)=\int_a^tD^r\partial_sS(s)\,\mathrm ds,
 \]
 故 \(D^rS\) 也有界。这个归纳从刚证出的 \(r=0\) 开始，给出全部固定背景联络导数。

 最后，度量相容式 \(\nabla g=0\) 给出
 \[
  (D_Xg)(Y,Z)=g(S(X,Y),Z)+g(Y,S(X,Z)).
 \]
 对它反复施 \(D\)，以已控制的 \(D^\ell S\) 及低阶 \(D^jg\) 归纳，得到所有 \(D^rg\) 的界。同样的联络转换把 \(\nabla^r\operatorname{Ric}\) 的界变为 \(D^r\operatorname{Ric}\) 的界。再由 \(\partial_tg=-2\operatorname{Ric}\)，得到所需的时间导数界。
\end{proof}

\begin{theorem}\label{g3:thm:ricci-flow-curvature-continuation}
 设 \(M^n\) 闭，\(n\ge2\)，\(g(t)\) 是 \([0,T)\)、\(0<T<\infty\)，上的光滑 Ricci 流。若
 \[
  \sup_{M\times[0,T)}|\operatorname{Rm}_{g(t)}|_{g(t)}<\infty,
 \]
 则存在 \(\varepsilon>0\)，使该流光滑延拓到 \([0,T+\varepsilon)\)。特别地，若最大存在时间 \(T_{\max}<\infty\)，则
 \[
  \limsup_{t\uparrow T_{\max}}
     \max_M|\operatorname{Rm}_{g(t)}|_{g(t)}=\infty.
 \]
\end{theorem}
\begin{proof}
 取 \(a=T/2\)。上一引理说明，对每个 \(r\)，\(g(t)\) 在固定背景的 \(C^r\) 范数中关于时间一致 Lipschitz，因为 \(\partial_tg\) 的该范数有界。故 \(t\uparrow T\) 时，\(g(t)\) 在每个 \(C^r\) 中收敛；这些极限由 \(C^0\) 极限唯一识别，组成一个光滑张量 \(g_T\)。式\eqref{g3:eq:ricci-metric-equivalence}保证 \(g_T\) 仍为正定度量。

 以 \(g_T\) 为初值应用\cref{g3:thm:ricci-flow-short-time}，得到 \(T\) 之后的光滑 Ricci 流。接合处的全部空间导数相同；一阶时间导数由 \(-2\operatorname{Ric}\) 决定，反复微分流方程使每个更高时间导数成为度量的有限阶空间喷射的光滑表达。两侧的这些喷射都趋向同一个 \(g_T\) 的喷射，所以全部混合导数相接，得到光滑延拓。

 若最大有限时间处的曲率上极限有限，取接近该时刻的一段区间便有统一曲率界，之前的紧时间区间也因光滑性而有界。刚证出的延拓于是与最大性矛盾，得到末项。
\end{proof}

延拓判据把有限时间的失效定位为曲率失控。它的证明中，曲率本身先保证度量正定性不退化；协变导数估计控制几何的高阶变化，联络差再把这些控制转成固定坐标中的光滑极限。下一节完成其中使用的演化计算和导数估计。
% 纲要条目排为编号小节标题，并进入目录与书签；说明段仍保留于纲要注释中。

% !TeX root = ../../Book3.tex
% 纲要标题：### 27.3 曲率演化和保持性质
% 纲要位置：几何学计划纲要.md，第 3258 行。
% 正文状态：初稿；尚未正式审核。

% 纲要内容开始
% 
% * 体积、标量及 Ricci 演化
% * 完整曲率张量的反应项
% * 标量最大值与张量最大值的区别
% * 所选曲率锥的保持
% 
% 纲要内容结束

\section{曲率演化和保持性质}
\label{g3:sec:2703}

Ricci 流首先改变度量，而曲率是度量的二阶量。把曲率自身的演化写出来，可以直接观察扩散与二次反应的竞争，并检验一个曲率条件是否随时间保持。标量曲率只记录曲率张量的一个收缩；讨论所有二维方向的曲率，还须保留完整的四阶张量。本节先计算这些量，再用标量最大值估计曲率导数，最后对曲率算子的非负锥证明一个真正的张量保持结论。

\subsection{体积、标量及 Ricci 演化}
\label{g3:subsec:2703:01}
本节令 \(g(t)\) 为光滑 Ricci 流，\(\partial_tg=-2\operatorname{Ric}\)，并写 \(s=\operatorname{Scal}_{g(t)}\)。所有协变导数、内积和收缩都采用当前度量。曲率槽位沿用\cref{g3:def:curvature-four-tensor}：
\[
	R_{ijkl}=g(R(e_i,e_j)e_l,e_k),
	\qquad \operatorname{Ric}_{jl}=g^{ik}R_{ijkl}.
\]
所以球面的 \(R_{1212}\) 为正。对任意张量，\(\Delta=\nabla^*\nabla=-\operatorname{tr}_g\nabla^2\) 为非负粗 Laplace 算子；对函数它就是此前的非负 Laplace 算子。

\begin{proposition}\label{g3:prop:ricci-flow-basic-evolution}
	设 \(g(t)\) 是 \(n\) 维光滑流形上的光滑 Ricci 流，\(n\ge2\)。则
	\begin{align}
		\partial_t\mu_g&=-s\mu_g,\label{g3:eq:ricci-flow-volume-evolution}\\
		(\partial_t+\Delta)\operatorname{Ric}_{ij}
		&=2\mathring R(\operatorname{Ric})_{ij}
		-2\operatorname{Ric}_i{}^a\operatorname{Ric}_{aj},\label{g3:eq:ricci-flow-ricci-evolution}\\
		(\partial_t+\Delta)s&=2|\operatorname{Ric}|^2.\label{g3:eq:ricci-flow-scalar-evolution}
	\end{align}
	其中曲率作用采用式\eqref{g3:eq:curvature-action-symmetric-tensors}，在正交标架中为
	\[
		\mathring R(\operatorname{Ric})_{ij}
		=\sum_{a,b}R_{a i b j}\operatorname{Ric}_{ab}.
	\]
	若流形闭，则
	\[
		\frac{\mathrm d}{\mathrm dt}\operatorname{Vol}_{g(t)}M
		=-\int_M s\,\mu_g.
	\]
\end{proposition}
\begin{proof}
	由式\eqref{g3:eq:inverse-volume-variation}，体积密度的导数是 \(\tfrac12\operatorname{tr}_g(\partial_tg)\mu_g\)，即第一式。收缩 Bianchi 恒等式给出 \(\delta\operatorname{Ric}+\tfrac12\mathrm ds=0\)，所以 \(\beta_g\operatorname{Ric}=0\)。在\cref{g3:thm:ricci-scalar-linearization}中令 \(h=-2\operatorname{Ric}\)，便有
	\[
		\partial_t\operatorname{Ric}
		=-\Delta_L\operatorname{Ric}
		=-\Delta\operatorname{Ric}
		-2\operatorname{Ric}\circ\operatorname{Ric}
		+2\mathring R(\operatorname{Ric}).
	\]
	这证明第二式。求 \(s=g^{ij}\operatorname{Ric}_{ij}\) 的导数时，\(\partial_tg^{ij}=2\operatorname{Ric}^{ij}\) 额外贡献 \(2|\operatorname{Ric}|^2\)。另一方面，\(\operatorname{tr}\mathring R(\operatorname{Ric})=|\operatorname{Ric}|^2\)，故第二式右端取迹为零。因此得到第三式。闭性允许在体积积分中交换求导与积分，给出末项。
\end{proof}

\begin{proposition}\label{g3:prop:ricci-flow-scalar-lower-bound}
	设 \(M^n\) 闭，\(n\ge2\)，\(g(t)\)、\(0\le t<T\)，为光滑 Ricci 流。若 \(s(0)\ge a\)，则在分母为正的共同时间内
	\begin{equation}\label{g3:eq:ricci-flow-scalar-lower-bound}
		s(t)\ge\frac{a}{1-2at/n}.
	\end{equation}
	特别地，非负标量曲率保持；若 \(a>0\)，光滑存在时间满足 \(T\le n/(2a)\)。
\end{proposition}
\begin{proof}
	Cauchy 不等式给 \(|\operatorname{Ric}|^2\ge s^2/n\)。令 \(b(t)=a/(1-2at/n)\)，则 \(b'=2b^2/n\)。对 \(w=b-s\)，式\eqref{g3:eq:ricci-flow-scalar-evolution}给出
	\[
		(\partial_t+\Delta)w
		\le\frac2n(b^2-s^2)=\frac2n(b+s)w,
		\qquad w(0)\le0.
	\]
	在任意所述范围内的紧时间区间上，系数有界，\(g(t)\) 一致正定。将其 Laplace 算子用一个固定背景联络表达，便符合\cref{g3:thm:parabolic-maximum}的时间依赖系数条件，因此 \(w\le0\)。若 \(a>0\) 且流光滑存在超过 \(n/(2a)\)，右端在一个内时刻前发散，与光滑曲率有界矛盾。
\end{proof}

标量曲率非负时，总体积非增；若标量曲率有负区域，体积的变化取决于它的积分。即使 \(s\) 已经有下界，曲率在不同方向间的抵消仍可能使完整张量很大，因而延拓时需要进一步控制 \(\mathrm{Rm}\)。

\subsection{完整曲率张量的反应项}
\label{g3:subsec:2703:02}
四阶张量的演化包含两类二次项。一类来自交换协变导数，另一类来自度量和指标长度的变化。先在固定坐标中写出全部项，再使用随时间保持正交的标架消去后一类项。

\begin{proposition}\label{g3:prop:ricci-flow-full-curvature}
	设 \(g(t)\) 是 \(n\) 维光滑 Ricci 流，\(n\ge2\)。定义二次曲率张量
	\begin{equation}\label{g3:eq:ricci-flow-quadratic-curvature}
		Q_{ijkl}
		=\sum_{p,q}R_{ijpq}R_{klpq}
		+2\sum_{p,q}\bigl(R_{ipkq}R_{jplq}-R_{iplq}R_{jpkq}\bigr),
	\end{equation}
	其中在当前正交标架中求和；一般坐标下在每对求和指标间插入逆度量。固定源坐标中的协变四阶张量满足
	\begin{equation}\label{g3:eq:ricci-flow-curvature-coordinate-evolution}
		\begin{split}
		(\partial_t+\Delta)R_{ijkl}={}&Q_{ijkl}
		-\operatorname{Ric}_i{}^pR_{pjkl}
		-\operatorname{Ric}_j{}^pR_{ipkl}\\
		&-\operatorname{Ric}_k{}^pR_{ijpl}
		-\operatorname{Ric}_l{}^pR_{ijkp}.
		\end{split}
	\end{equation}
\end{proposition}
\begin{proof}
	记 \(r_{ij}=\operatorname{Ric}_{ij}\)。联络变分公式\eqref{g3:eq:connection-variation}给
	\[
		C^a_{ij}=\partial_t\Gamma^a_{ij}
		=-g^{ab}(\nabla_i r_{jb}+\nabla_j r_{ib}-\nabla_b r_{ij}).
	\]
	在曲率坐标式中微分，得到
	\(\partial_tR^a{}_{l i j}=\nabla_i C^a_{jl}-\nabla_j C^a_{il}\)。再微分降指标所用的 \(g_{ka}\)，并交换作用于 \(r_{kl}\) 的两次导数，可得
	\begin{equation}\label{g3:eq:ricci-flow-curvature-uncommuted}
		\partial_tR_{ijkl}
		=D_{ijkl}-r_k{}^pR_{ijpl}-r_l{}^pR_{ijkp},
	\end{equation}
	其中
	\[
		D_{ijkl}=\nabla_i\nabla_k r_{jl}-\nabla_i\nabla_l r_{jk}
		-\nabla_j\nabla_k r_{il}+\nabla_j\nabla_l r_{ik}.
	\]
	为核对这次交换，尚未交换时另有
	\((\nabla_j\nabla_i-\nabla_i\nabla_j)r_{lk}-2r_k{}^pR_{ijpl}\)。前一项等于 \(r_k{}^pR_{ijpl}-r_l{}^pR_{ijkp}\)，故其和恰为式\eqref{g3:eq:ricci-flow-curvature-uncommuted}中的两项。

	下面把 \(D\) 改写成 Laplace 项。第二 Bianchi 恒等式及其收缩给出
	\[
		\nabla_pR_{ijkl}=\nabla_iR_{pjkl}-\nabla_jR_{pikl},
		\qquad \nabla^pR_{pjkl}=\nabla_k r_{jl}-\nabla_l r_{jk}.
	\]
	在所算点取正交法标架，再求一次导数，便有
	\begin{equation}\label{g3:eq:ricci-flow-curvature-laplace-commutator}
		-\Delta R_{ijkl}=D_{ijkl}
		+\sum_p\bigl([\nabla_p,\nabla_i]R\bigr)_{pjkl}
		-\sum_p\bigl([\nabla_p,\nabla_j]R\bigr)_{pikl}.
	\end{equation}
	四个协变槽位的交换公式是
	\[
		([\nabla_p,\nabla_i]R)_{abcd}
		=-R_{piua}R_{ubcd}-R_{piub}R_{aucd}
		-R_{piuc}R_{abud}-R_{piud}R_{abcu}.
	\]
	在式\eqref{g3:eq:ricci-flow-curvature-laplace-commutator}中，作用于第一个槽位的两项收缩为
	\(r_i{}^pR_{pjkl}+r_j{}^pR_{ipkl}\)。作用于第二个槽位的两项，由第一 Bianchi 恒等式
	\(R_{pjui}-R_{piuj}=-R_{ijpu}\)，合为
	\(-\sum_{p,u}R_{ijpu}R_{klpu}\)。余下两个槽位的四项，利用成对反对称性合为
	\[
		-2\sum_{p,u}R_{ipku}R_{jplu}
		+2\sum_{p,u}R_{iplu}R_{jpku}.
	\]
	因此全部交换项为 \(r_i{}^pR_{pjkl}+r_j{}^pR_{ipkl}-Q_{ijkl}\)。解出 \(D\) 并代回式\eqref{g3:eq:ricci-flow-curvature-uncommuted}，即得所求完整演化。
\end{proof}

为了比较特征值，把变化的内积搬到一个固定的欧氏向量丛上。令 \(E=\mathrm TM\)，但固定其内积为 \(h=g(0)\)。逐点解线性常微分方程
\begin{equation}\label{g3:eq:ricci-flow-orthonormal-bundle-ode}
	\partial_tP_t=\operatorname{Ric}_{g(t)}^\sharp\circ P_t,
	\qquad P_0=\operatorname{Id},
\end{equation}
得到丛同构 \(P_t:E\to\mathrm TM\)。对固定截面 \(v,w\)，直接求导得
\[
	\partial_t\bigl(g(t)(P_tv,P_tw)\bigr)
	=-2\operatorname{Ric}(P_tv,P_tw)
	+2\operatorname{Ric}(P_tv,P_tw)=0.
\]
所以 \(P_t\) 是从 \((E,h)\) 到 \((\mathrm TM,g(t))\) 的等距同构。定义拉回联络
\(D^t_Xv=P_t^{-1}\nabla^{g(t)}_X(P_tv)\)，它保持固定内积 \(h\)。以下 \(\widehat\Delta\) 是此联络与基底度量 \(g(t)\) 所定义的非负粗 Laplace 算子。

\ProseParagraphBreak
将四阶曲率沿四个输入都用 \(P_t\) 拉回，得到 \(\widehat R\)，并以固定 \(h\) 把它看成 \(\Lambda^2E\) 上的自伴算子 \(\widehat{\mathcal R}\)。对 \(\widehat R\) 求时间导数时，四个输入的导数恰抵消式\eqref{g3:eq:ricci-flow-curvature-coordinate-evolution}中的四个 Ricci 项。因此
\[
	(\partial_t+\widehat\Delta)\widehat R=Q(\widehat R).
\]
这个等式的二次项还可用二向量的代数表示，非负性由此变得可见。

\begin{definition}\label{g3:def:ricci-flow-lie-square}
	设 \((V,h)\) 为有限维欧氏空间。通过
	\[
		\iota(u\wedge v)z=h(v,z)u-h(u,z)v
	\]
	把 \(\Lambda^2V\) 与反对称自同态空间 \(\mathfrak{so}(V)\) 等距识别，后者内积为 \(-\tfrac12\operatorname{tr}(AB)\)。在 \(\Lambda^2V\) 上用矩阵交换子定义括号。取任意正交单位基 \((E_\alpha)\)，记
	\(c_{\alpha\beta\gamma}=\langle[E_\alpha,E_\beta],E_\gamma\rangle\)。对自伴算子 \(\mathcal R\)，定义自伴算子 \(\mathcal R^\#\) 的分量为
	\begin{equation}\label{g3:eq:ricci-flow-lie-square}
		(\mathcal R^\#)_{\alpha\beta}
		=\frac12\sum_{\gamma,\delta,\eta,\theta}
		c_{\alpha\gamma\delta}c_{\beta\eta\theta}
		\mathcal R_{\gamma\eta}\mathcal R_{\delta\theta}.
	\end{equation}
	所有指标都是同一欧氏空间上的完全收缩，故此定义与正交基的选择无关。
\end{definition}

\begin{proposition}\label{g3:prop:ricci-flow-operator-evolution}
	上述固定欧氏丛上的曲率算子满足
	\begin{equation}\label{g3:eq:ricci-flow-operator-evolution}
		(\partial_t+\widehat\Delta)\widehat{\mathcal R}
		=2\bigl(\widehat{\mathcal R}^{2}
		+\widehat{\mathcal R}^{\#}\bigr).
	\end{equation}
	归一化采用 \(e_i\wedge e_j\)、\(i<j\)，为正交单位基；在这一约定下常截面曲率 \(\kappa\) 对应 \(\mathcal R=\kappa\operatorname{Id}\)。
\end{proposition}
\begin{proof}
	暂略去帽号。在 \(E_{ij}=e_i\wedge e_j\) 中允许任意指标顺序，并规定 \(E_{ji}=-E_{ij}\)、\(E_{ii}=0\)。矩阵乘法直接给出
	\[
		[E_{ij},E_{kl}]
		=\delta_{jk}E_{il}-\delta_{ik}E_{jl}
		-\delta_{jl}E_{ik}+\delta_{il}E_{jk}.
	\]
	将这个括号代入式\eqref{g3:eq:ricci-flow-lie-square}，并把有序二指标求和改成对全部指标求和，得到
	\[
		\langle\mathcal R^\# E_{ij},E_{kl}\rangle
		=\sum_{p,q}\bigl(R_{ipkq}R_{jplq}-R_{iplq}R_{jpkq}\bigr).
	\]
	这里括号中的四个 Kronecker 项依次收缩公共指标；同一反对称二向量的两个次序给出相同乘积，正由定义中的 \(1/2\) 消去重复。另一方面，算子平方为
	\[
		\langle\mathcal R^2 E_{ij},E_{kl}\rangle
		=\sum_{p<q}R_{ijpq}R_{klpq}
		=\frac12\sum_{p,q}R_{ijpq}R_{klpq}.
	\]
	比较式\eqref{g3:eq:ricci-flow-quadratic-curvature}，可见 \(Q=2(\mathcal R^2+\mathcal R^\#)\)。结合拉回张量的演化即得结论。
\end{proof}

\begin{example}\label{g3:exm:ricci-flow-curvature-factor-check}
	若初始度量的常截面曲率为 \(\kappa_0\)，则 \(g(t)=(1-2(n-1)\kappa_0t)g(0)\)，而
	\(\kappa(t)=\kappa_0/(1-2(n-1)\kappa_0t)\)。对 \(\mathcal R=\kappa\operatorname{Id}\)，前述括号公式给
	\[
		\mathcal R^2=\kappa^2\operatorname{Id},\qquad
		\mathcal R^\#=(n-2)\kappa^2\operatorname{Id}.
	\]
	例如 \(E_{12}\) 可由每个 \(p\ne1,2\) 对应的一对 \(E_{1p},E_{p2}\) 的括号得到；交换这两个输入再给一次相同的平方，\(1/2\) 留下 \(n-2\) 个贡献。因此式\eqref{g3:eq:ricci-flow-operator-evolution}给出 \(\kappa'=2(n-1)\kappa^2\)，与显式度量完全一致。在二维 \(\mathfrak{so}(2)\) 的括号为零，\(\mathcal R^\#=0\)，但三维及以上通常必须保留这一项。
\end{example}

\subsection{标量最大值与张量最大值的区别}
\label{g3:subsec:2703:03}
标量函数在极小点具有非负 Hessian，因而有\cref{g3:thm:parabolic-maximum}的比较原理。张量的某个坐标分量为负却没有坐标不变的含义；曲率算子非负，要求它对每个二向量的二次型都非负。证明这种保持性，必须检验二次反应项在锥的边界朝哪个方向。此前可先把张量转为长度平方，用标量比较得到延拓需要的导数界。

\ProseParagraphBreak
以下用 \(U*V\) 表示由 \(g\)、\(g^{-1}\) 对 \(U\otimes V\) 作有限次缩并并线性组合。星号不指定符号，但每一项的范数有 \(|U*V|\le C|U||V|\)，常数只依赖维数、张量阶数和所列组合。涉及保持性的精确反应项仍使用式\eqref{g3:eq:ricci-flow-quadratic-curvature}；下面的星号只用于导数的长度估计。

\begin{lemma}\label{g3:lem:ricci-flow-derivative-evolution}
	设 \(g(t)\) 是光滑 \(n\) 维 Ricci 流，\(n\ge2\)。对整数 \(k\ge0\)，记
	\(F_k=|\nabla^k\mathrm{Rm}|^2\)、\(P=\partial_t+\Delta\)。则张量方程中的反应项具有形式
	\begin{equation}\label{g3:eq:ricci-flow-derivative-types}
		(\partial_t+\Delta)\nabla^k\mathrm{Rm}
		=\sum_{a+b=k}\nabla^a\mathrm{Rm}*\nabla^b\mathrm{Rm},
	\end{equation}
	其中每个星号为一个确定的有限线性组合。因而
	\begin{equation}\label{g3:eq:ricci-flow-derivative-norm}
		PF_k\le-2F_{k+1}
		+C_{n,k}\sum_{a=0}^k
		|\nabla^a\mathrm{Rm}|\,
		|\nabla^{k-a}\mathrm{Rm}|\,|\nabla^k\mathrm{Rm}|.
	\end{equation}
	特别地，\(PF_0\le-2F_1+C_n|\mathrm{Rm}|^3\)。若 \(|\mathrm{Rm}|\) 及所有低于 \(k\) 阶的曲率导数有给定上界，则
	\begin{equation}\label{g3:eq:ricci-flow-derivative-inductive-bound}
		PF_k\le-2F_{k+1}+C_k(1+F_k),
	\end{equation}
	其中 \(C_k\) 只依赖维数、阶数及这些低阶界。
\end{lemma}
\begin{proof}
	\(k=0\) 的张量等式是\cref{g3:prop:ricci-flow-full-curvature}；其所有零阶项都为 \(\mathrm{Rm}*\mathrm{Rm}\)。说明求导时没有遗漏更高阶项：对任意协变张量 \(A\)，联络变分给出
	\[
		[\partial_t,\nabla]A=\nabla\operatorname{Ric}*A,
	\]
	因为每个协变槽位只产生一个 \(\partial_t\Gamma\) 项，而 \(\partial_t\Gamma=g^{-1}*\nabla\operatorname{Ric}\)。将两次空间导数交换再求一次导数，得到
	\[
		[\Delta,\nabla]A
		=\mathrm{Rm}*\nabla A+\nabla\mathrm{Rm}*A.
	\]
	这来自每个槽位的曲率作用及其一次协变导数；\(\nabla g=0\) 保证迹的系数不另生空间导数。假定式\eqref{g3:eq:ricci-flow-derivative-types}在 \(k-1\) 阶成立，对它求一次协变导数，乘积中的总导数阶增加一。新增交换项为
	\(\mathrm{Rm}*\nabla^k\mathrm{Rm}\) 或
	\(\nabla\mathrm{Rm}*\nabla^{k-1}\mathrm{Rm}\)，总阶同样为 \(k\)。这证明归纳公式。

	对 \(A=\nabla^k\mathrm{Rm}\) 求范数平方的时间导数。范数中的每个逆度量都贡献一个 \(\operatorname{Ric}*A*A\) 项，故其总绝对值至多为 \(C_{n,k}|\mathrm{Rm}||A|^2\)。空间乘积法则则给
	\[
		\Delta|A|^2=2\langle\Delta A,A\rangle-2|\nabla A|^2.
	\]
	结合张量方程，便得式\eqref{g3:eq:ricci-flow-derivative-norm}，度量变化项可并入 \(a=0\) 或 \(a=k\) 的项。对中间的 \(1\le a\le k-1\)，两个曲率因子都由低阶界控制；乘以 \(|A|\) 后用 \(|A|\le(1+|A|^2)/2\)。两端项直接为有界常数乘 \(|A|^2\)。这就得到最后的不等式。
\end{proof}

\begin{theorem}\label{g3:thm:ricci-flow-curvature-derivatives}
	设 \(M^n\) 闭，\(n\ge2\)，\(g(t)\)、\(0\le t<T\)，为光滑 Ricci 流，\(T\le\infty\)，并满足
	\[
		\sup_{M\times[0,T)}|\mathrm{Rm}|\le K<\infty.
	\]
	对每个整数 \(k\ge1\) 及 \(0<\varepsilon<T\)，存在只依赖 \(n,k,K,\varepsilon\) 的常数，使
	\begin{equation}\label{g3:eq:ricci-flow-positive-time-derivative-bound}
		\sup_{\varepsilon\le t<T}|\nabla^k\mathrm{Rm}|\le C(n,k,K,\varepsilon).
	\end{equation}
	常数不依赖初始曲率的高阶导数或终止时间 \(T\)。
\end{theorem}
\begin{proof}
	固定 \(k,\varepsilon\)，令 \(\delta=\varepsilon/k\)。逐阶证明 \(F_j\) 在 \(t\ge j\delta\) 上有统一界，\(0\le j\le k\)。\(j=0\) 是假设。设 \(F_0,\ldots,F_{j-1}\) 已在 \(t\ge(j-1)\delta\) 上有界，记其中 \(F_{j-1}\le B\)。任取 \(t\ge j\delta\)，令 \(a=t-\delta\)，在时间窗 \([a,t]\) 上考虑
	\[
		Z(x,\sigma)=(\sigma-a)F_j(x,\sigma)+A F_{j-1}(x,\sigma),
		\qquad a\le\sigma\le t.
	\]
	上一引理给出与窗口位置无关的常数 \(C,C'\)，使
	\[
		PF_j\le-2F_{j+1}+C(1+F_j),\qquad
		PF_{j-1}\le-2F_j+C'.
	\]
	后一式在 \(j=1\) 时用 \(PF_0\le-2F_1+C_nK^3\)；在 \(j>1\) 时使用已知的低阶界。因此
	\[
		PZ\le(1+\delta C-2A)F_j+\delta C+AC'.
	\]
	取 \(2A\ge2+\delta C\)，便有 \(PZ\le C''\)，其中常数依赖已列数据而与 \(a\) 无关。初值 \(Z(\cdot,a)=AF_{j-1}(\cdot,a)\le AB\) 不涉及 \(F_j(a)\)。标量最大值原理给出
	\[
		\delta F_j(x,t)\le Z(x,t)\le AB+C''\delta.
	\]
	因为 \(t\) 任意，这完成归纳。对 \(j=k\)，开始时间正是 \(\varepsilon\)。每个窗口上的度量虽在变化，但光滑正定性保证可用最大值原理；所得上界只使用右端和初值，并不引入随窗口位置变化的坐标椭圆常数。
\end{proof}

这个证明每次用低一阶量的负耗散项控制高一阶量。乘上窗口内的时间长度，使最高阶的初始值从 \(Z(a)\) 中消失；因此只要离初时刻一段正时间，就不需要初始高阶曲率界。\secref{g3:sec:2702}的有限时延拓正使用这一结论，把曲率控制提升为能够重新求解的光滑数据。

\subsection{所选曲率锥的保持}
\label{g3:subsec:2703:04}
现在取一个在所有维数都明确的条件：曲率算子作为 \(\Lambda^2\mathrm TM\) 上的自伴算子非负。它比只在可分解二向量上测试的非负截面曲率更直接地适配式\eqref{g3:eq:ricci-flow-operator-evolution}。在固定欧氏丛上，这些非负算子组成闭凸锥。

\begin{lemma}\label{g3:lem:ricci-flow-lie-square-positive}
	设 \(V\) 为有限维欧氏空间，\(\mathcal S:\Lambda^2V\to\Lambda^2V\) 自伴且非负。则 \(\mathcal S^2\) 与 \(\mathcal S^\#\) 都非负。特别地，若 \(\mathcal S\xi=0\)，则
	\[
		\left\langle2(\mathcal S^2+\mathcal S^\#)\xi,\xi\right\rangle\ge0.
	\]
\end{lemma}
\begin{proof}
	取 \(\mathcal S\) 的正交特征基 \((E_\gamma)\)，特征值 \(\lambda_\gamma\ge0\)。对任意 \(\xi\)，式\eqref{g3:eq:ricci-flow-lie-square}成为
	\[
		\langle\mathcal S^\#\xi,\xi\rangle
		=\frac12\sum_{\gamma,\delta}
		\lambda_\gamma\lambda_\delta
		\langle[\xi,E_\gamma],E_\delta\rangle^2\ge0.
	\]
	而 \(\langle\mathcal S^2\xi,\xi\rangle=|\mathcal S\xi|^2\ge0\)。这同时证明全部结论。
\end{proof}

\begin{theorem}\label{g3:thm:ricci-flow-nonnegative-curvature-operator}
	设 \(M^n\) 闭，\(n\ge2\)，\(g(t)\)、\(0\le t<T\)，为光滑 Ricci 流。若 \(g(0)\) 的曲率算子非负，则每个 \(g(t)\) 的曲率算子都非负。若初始曲率算子正定，则它在全部光滑存在时间内保持正定。
\end{theorem}
\begin{proof}
	在式\eqref{g3:eq:ricci-flow-orthonormal-bundle-ode}给出的固定内积丛上工作，略去帽号，并写
	\(\mathcal Q(\mathcal R)=2(\mathcal R^2+\mathcal R^\#)\)、\(L=-\widehat\Delta\)。固定 \(T_1<T\)，令 \(K_1=\sup_{M\times[0,T_1]}|\mathcal R|\)。由于 \(\mathcal Q\) 是有限维二次多项式，存在只依赖 \(n\) 的 \(C_n\)，使 \(0\le b\le1\) 时
	\[
		|\mathcal Q(\mathcal R+bI)-\mathcal Q(\mathcal R)|
		\le C_n b(K_1+1).
	\]
	选 \(C>C_n(K_1+1)\)，再取足够小的 \(\eta>0\)，使 \(b(t)=\eta e^{Ct}\le1\) 于 \([0,T_1]\)。令 \(\mathcal S=\mathcal R+bI\)。初始 \(\mathcal S\) 正定。

	若 \(\mathcal S\) 第一次在 \((x_*,t_*)\) 出现零特征值，取单位零向量 \(\xi\)。此时 \(\mathcal S\ge0\)，所以 \(\mathcal S\xi=0\)。在该点附近把 \(\xi\) 延拓成局部截面，使其一阶空间协变导数在所算点为零。标量 \(\langle\mathcal S\xi,\xi\rangle\) 此前为正、此时有零的空间最小值；其时间导数非正、\(L\) 非负。计算中的额外导数项含 \(D\xi\) 或 \(\mathcal S\xi\)，都在该点消失。故
	\[
		0\ge\langle(\partial_t-L)\mathcal S\xi,\xi\rangle
		=\langle\mathcal Q(\mathcal R)\xi,\xi\rangle+Cb.
	\]
	另一方面，上一引理给 \(\langle\mathcal Q(\mathcal S)\xi,\xi\rangle\ge0\)，从而右端至少为
	\(b\bigl(C-C_n(K_1+1)\bigr)>0\)，矛盾。因此 \(\mathcal R+\eta e^{Ct}I\) 始终正定；令 \(\eta\downarrow0\)，得到非负性。

	若初值正定，紧性允许取 \(b_0>0\)，使 \(\mathcal R(0)>b_0I\)。这次令 \(b(t)=b_0e^{-t}\)、\(\mathcal S=\mathcal R-bI\)。在假设的第一个零方向上，\(\mathcal R\ge bI\ge0\)，所以 \(\langle\mathcal Q(\mathcal R)\xi,\xi\rangle\ge0\)。相同的最小值论证却给出
	\[
		0\ge\langle(\partial_t-L)\mathcal S\xi,\xi\rangle
		=\langle\mathcal Q(\mathcal R)\xi,\xi\rangle-b'
		\ge b>0.
	\]
	故 \(\mathcal R>b_0e^{-t}I\)，特别地始终正定。\(T_1<T\) 任意，证明完成。
\end{proof}

这里真正决定保持性的，是 \(\mathcal Q\) 在非负锥边界的方向以及连接保持纤维内积。证明没有把曲率分量逐个套入标量最大值原理。所得条件还立即蕴含非负截面曲率、非负 Ricci 曲率及非负标量曲率：依次在可分解单位二向量上取值，再沿正交方向求和即可。

\begin{example}\label{g3:exm:ricci-flow-product-curvature-cone}
	取 \(p\ge2\)、\(q\ge1\)，在 \(S^p\times\mathbb T^q\) 上使用单位圆球与平坦环面的乘积度量。乘积联络不混合两个因子，曲率算子在 \(\Lambda^2\mathrm TS^p\) 上为恒等算子，在混合二向量和 \(\Lambda^2\mathrm T\mathbb T^q\) 上为零。Ricci 流为
	\[
		g(t)=(1-2(p-1)t)g_{S^p}\oplus g_{\mathbb T^q},
		\qquad 0\le t<\frac1{2(p-1)}.
	\]
	其正特征值为 \((1-2(p-1)t)^{-1}\)，其余特征值始终为零。因此非负锥保持，并不表示边界上的所有零特征值立即变为正值。这个例子也保留了高维混合方向，不能由标量曲率的一个方程单独描述。
\end{example}
% 纲要条目排为编号小节标题，并进入目录与书签；说明段仍保留于纲要注释中。

% !TeX root = ../../Book3.tex
% 纲要标题：### 27.4 放大、孤立子和熵
% 纲要位置：几何学计划纲要.md，第 3265 行。
% 正文状态：初稿，尚未正式审核。

% 纲要内容开始
% 
% * 奇点尺度及抛物重缩放
% * 收缩、稳态和膨胀孤立子
% * Perelman 单调泛函
% * 有限时间非塌缩
% 
% 纲要内容结束

\section{放大、孤立子和熵}
\label{g3:sec:2704}

曲率变大时，原来的长度单位已不适合观察局部形状。先把曲率所在的尺度放大，再寻找在缩放和微分同胚作用下保持形状的解，便得到奇点分析的两个基本对象：重缩放流和孤立子。单调泛函进一步辨认何时演化只剩下这种自相似运动；非塌缩定理则给放大后的球提供体积下界。

\subsection{奇点尺度及抛物重缩放}
\label{g3:subsec:2704:01}
设 \(M^n\) 为光滑流形，\(n\ge2\)，\(g(t)\) 满足 \(\partial_tg=-2\operatorname{Ric}_g\)。若 \(a>0\) 为常数，则 \(ag\) 与 \(g\) 的 Levi--Civita 联络相同，因而作为 \((0,2)\) 型张量的 Ricci 曲率不变，而
\begin{equation}\label{g3:eq:ricci-constant-scaling}
	|\operatorname{Rm}_{ag}|_{ag}=a^{-1}|\operatorname{Rm}_g|_g,
	\qquad s_{ag}=a^{-1}s_g,
	\qquad \mu_{ag}=a^{n/2}\mu_g.
\end{equation}
第一式也可直接按指标核对：全协变曲率张量乘以 \(a\)，其平方范数中的四个逆度量各乘以 \(a^{-1}\)。长度则乘以 \(a^{1/2}\)。

\ProseParagraphBreak
给定时空点 \((x_j,t_j)\) 及 \(Q_j>0\)，定义\term[paowu chongsuofang]{抛物重缩放}[parabolic rescaling]
\begin{equation}\label{g3:eq:ricci-parabolic-rescaling}
	g_j(\sigma)\coloneqq Q_jg(t_j+Q_j^{-1}\sigma).
\end{equation}
链式法则和 Ricci 张量的缩放性质给出
\[
	\partial_\sigma g_j=-2\operatorname{Ric}_{g_j}.
\]
这说明空间长度放大 \(Q_j^{1/2}\) 时，时间长度必须放大 \(Q_j\)。若原流定义于 \([0,T)\)，则新时间区间为
\[
	[-Q_jt_j,\ Q_j(T-t_j)).
\]
取 \(Q_j=|\operatorname{Rm}_{g(t_j)}|(x_j)>0\)，便有 \(|\operatorname{Rm}_{g_j(0)}|(x_j)=1\)；原来的曲率尺度 \(Q_j^{-1/2}\) 被变成单位长度。

\begin{definition}\label{g3:def:ricci-singularity-types}
	设闭流形上的最大光滑 Ricci 流存在于 \([0,T)\)，且 \(T<\infty\)。记
	\[
		K(t)=\max_M|\operatorname{Rm}_{g(t)}|.
	\]
	若 \(\sup_{0\le t<T}(T-t)K(t)<\infty\)，称该有限时奇点为\term[I xing qidian]{I 型奇点}[type I singularity]；若此上确界为无穷，则称为\term[II xing qidian]{II 型奇点}[type II singularity]。
\end{definition}

这一区分记录曲率相对于 \((T-t)^{-1}\) 的增长速度。圆球收缩给出 I 型奇点：若初始半径为 \(r_0\)，则半径平方为 \(r_0^2-2(n-1)t\)，曲率范数与 \((T-t)^{-1}\) 成正比。定义本身没有断言奇点经放大后趋向某个固定模型。

\ProseParagraphBreak
\Cref{g3:thm:ricci-flow-curvature-continuation}保证最大解的有限时奇点处曲率无界。可取 \(\tau_j\uparrow T\)，再在 \(M\times[0,\tau_j]\) 选取曲率范数的最大点 \((x_j,t_j)\)，并记其值为 \(Q_j\)，使 \(Q_j\to\infty\)。每个较短的闭时间区间上曲率有界，故 \(t_j\to T\)。于是
\begin{equation}\label{g3:eq:ricci-backward-rescaled-bound}
	|\operatorname{Rm}_{g_j(\sigma)}|\le1
	\quad(-Q_jt_j\le\sigma\le0),
	\qquad |\operatorname{Rm}_{g_j(0)}|(x_j)=1.
\end{equation}
这里 \(Q_jt_j\to\infty\)，所以任意固定的向后时间区间最终都包含在定义域内。这已经把原来的无界曲率问题化为统一有界的重缩放序列。要从序列取得光滑极限，还须控制局部几何的退化并使用相应的导数估计和紧性定理；式\eqref{g3:eq:ricci-backward-rescaled-bound}只提供其中的曲率界。

\subsection{收缩、稳态和膨胀孤立子}
\label{g3:subsec:2704:02}
若一个流由固定度量通过随时间变化的常数缩放与微分同胚拉回得到，就称它为\term[zixiangsi jie]{自相似解}[self-similar solution]。它的形状在这些等价关系下保持不变。梯度型的情形可以用一个定常张量方程表达。

\begin{definition}\label{g3:def:gradient-ricci-soliton}
	设 \((M,g)\) 为光滑黎曼流形，\(f\in C^\infty(M)\)，\(\lambda\in\mathbb R\) 为常数。若
	\begin{equation}\label{g3:eq:gradient-ricci-soliton}
		\operatorname{Ric}_g+\operatorname{Hess}_g f=\lambda g,
	\end{equation}
	就称 \((g,f,\lambda)\) 为\term[tidu Ricci gulizi]{梯度 Ricci 孤立子}[gradient Ricci soliton]，\(f\) 为势函数。按 \(\lambda>0\)、\(\lambda=0\)、\(\lambda<0\)，分别称为\term[shousuo gulizi]{收缩孤立子}[shrinking soliton]、\term[wentai gulizi]{稳态孤立子}[steady soliton]和\term[pengzhang gulizi]{膨胀孤立子}[expanding soliton]。
\end{definition}

\begin{proposition}\label{g3:prop:soliton-selfsimilar-flow}
	设 \((M,g)\) 为光滑无边界黎曼流形，\(f\in C^\infty(M)\)、\(\lambda\in\mathbb R\)，且 \(\operatorname{Ric}_g+\operatorname{Hess}_gf=\lambda g\)，令 \(c(t)=1-2\lambda t\)。在包含 \(0\)、满足 \(c(t)>0\) 的时间区间上，设微分同胚族 \(\phi_t\) 解
	\[
		\frac{\mathrm d}{\mathrm dt}\phi_t
		=\frac1{c(t)}\nabla_g f\circ\phi_t,
		\qquad \phi_0=\operatorname{id}.
	\]
	则 \(g(t)=c(t)\phi_t^*g\) 是初值为 \(g\) 的 Ricci 流。若 \(M\) 闭，这一微分同胚族在 \(c>0\) 的整个连通时间区间上存在。
\end{proposition}
\begin{proof}
	利用 \(\mathcal L_{\nabla f}g=2\operatorname{Hess}f\)，得到
	\[
		\begin{aligned}
		\partial_t\bigl(c\phi_t^*g\bigr)
		&=\phi_t^*\bigl(-2\lambda g+2\operatorname{Hess}f\bigr)\\
		&=-2\phi_t^*\operatorname{Ric}_g
		=-2\operatorname{Ric}_{c\phi_t^*g}.
		\end{aligned}
	\]
	最后一步用了 Ricci 张量的自然性和常数缩放不变性。闭流形上的光滑向量场在每个紧时间区间上生成全局微分同胚；这些区间覆盖 \(c>0\) 的相应连通区间。非紧情形则需另行保证所用向量场的流存在。
\end{proof}

\begin{example}\label{g3:exm:gaussian-ricci-soliton}
	在 \(\mathbb R^n\) 上取 Euclid 度量 \(g_{\mathrm E}\) ，令 \(\lambda\in\mathbb R\)、\(c(t)=1-2\lambda t>0\)，取 \(f(x)=\lambda|x|^2/2\)。由于 \(\operatorname{Ric}=0\)、\(\operatorname{Hess}f=\lambda g_{\mathrm E}\)，这给出\term[Gauss Ricci gulizi]{Gauss 孤立子}[Gaussian soliton]。向量场为 \(\lambda x/c(t)\)，故 \(\phi_t(x)=c(t)^{-1/2}x\)。这里
	\[
		c(t)\phi_t^*g_{\mathrm E}=g_{\mathrm E}.
	\]
	因此同一个平坦静止流可以用收缩或膨胀的孤立子数据表示；分类名称说的是孤立子方程中的 \(\lambda\)，而非固定坐标下每一点的度量都必须严格缩小或增大。
\end{example}

\begin{example}\label{g3:exm:round-cylinder-soliton}
	半径 \(r\) 的圆球 \(S^n(r)\)，\(n\ge2\)，满足 \(\operatorname{Ric}=(n-1)r^{-2}g\)。取常势函数和 \(\lambda=(n-1)r^{-2}\)，便得到圆球收缩流。
	在乘积 \(S^k(r)\times\mathbb R^{n-k}\) 上，\(2\le k<n\)，取
	\[
		g=g_{S^k(r)}+g_{\mathrm E},\qquad
		\lambda=\frac{k-1}{r^2},\qquad
		f(y)=\frac\lambda2|y|^2.
	\]
	球面方向的 Ricci 张量等于 \(\lambda g\)，Euclid 方向的 Hessian 等于 \(\lambda g\)，混合项均为零。因此这是收缩孤立子。令 \(c(t)=1-2\lambda t>0\)，相应的 Ricci 流为
	\[
		g(t)=c(t)g_{S^k(r)}+g_{\mathrm E},
	\]
	它的球面因子收缩，而平直因子保持不变。
\end{example}

\begin{example}\label{g3:exm:cigar-ricci-soliton}
	在 \(\mathbb R^2\) 上令
	\[
		g=\frac{\mathrm dx^2+\mathrm dy^2}{1+x^2+y^2},
		\qquad f=-\log(1+x^2+y^2).
	\]
	这是 Hamilton 的\term[xuejia gulizi]{雪茄孤立子}[cigar soliton]。写 \(r=\sqrt{x^2+y^2}\)、\(\rho=\operatorname{arsinh}r\)，则
	\[
		g=\mathrm d\rho^2+\tanh^2\rho\,\mathrm d\theta^2,
		\qquad f=-2\log\cosh\rho.
	\]
	对 \(a(\rho)=\tanh\rho\)，Gauss 曲率为
	\[
		K=-\frac{a''}{a}=2\operatorname{sech}^2\rho.
	\]
	径向单位向量和角向单位向量上的 Hessian 分别为
	\[
		f''=-2\operatorname{sech}^2\rho,
		\qquad f'\frac{a'}a=-2\operatorname{sech}^2\rho,
	\]
	混合分量为零。二维中 \(\operatorname{Ric}=Kg\)，所以 \(\operatorname{Ric}+\operatorname{Hess}f=0\)。这些等式由光滑性延伸到原点。径向长度 \(\rho\) 趋向无穷，故度量完备；其曲率有界且为正。又因 \(\nabla_gf=-2r\partial_r\)，微分同胚是 \(\phi_t(x,y)=e^{-2t}(x,y)\)，于是得到对所有实数 \(t\) 定义的稳态孤立子流
	\[
		g(t)=\frac{\mathrm dx^2+\mathrm dy^2}{e^{4t}+x^2+y^2}.
	\]
\end{example}

\subsection{Perelman 单调泛函}
\label{g3:subsec:2704:03}
下面构造 Perelman 的两个单调泛函，并在闭流形上证明其导数公式。\footnote{\OriginalName{Григорий Яковлевич Перельман}（Grigori Yakovlevich Perelman，1966--），俄罗斯数学家，研究黎曼几何与几何分析，完成了三维几何化猜想的证明；人物资料见\href{https://mathshistory.st-andrews.ac.uk/Biographies/Perelman/}{MacTutor 数学史档案}。}本小节中 \(s=s_g\) 表示标量曲率，所有 Laplace 算子均取 \(\Delta=-\operatorname{div}\nabla\)。

\ProseParagraphBreak
对闭黎曼流形 \((M^n,g)\) 和光滑实函数 \(f\)，定义\term[Perelman F fanhan]{Perelman \(\mathcal F\) 泛函}[Perelman's \(\mathcal F\)-functional]
\begin{equation}\label{g3:eq:perelman-F-definition}
	\mathcal F(g,f)\coloneqq\int_M\bigl(s_g+|\mathrm df|_g^2\bigr)e^{-f}\mu_g.
\end{equation}
\symidx{\(\mathcal F(g,f)\)：Perelman 的 \(\mathcal F\) 泛函}
通常另取归一化条件 \(\int_Me^{-f}\mu_g=1\)。先不限制变分，以便看清权重与度量怎样相互配合。

\begin{lemma}\label{g3:lem:perelman-F-variation}
	设 \(M\) 为闭光滑流形，\(g(\varepsilon)\) 为光滑黎曼度量变分，\(f(\varepsilon)\) 为实光滑函数变分；记 \(\mathcal F(g,f)=\int_M(s_g+|\mathrm df|_g^2)e^{-f}\,\mathrm d\mu_g\)，\(\Delta_g=-\operatorname{div}_g\nabla\)，记 \(h=\dot g\)、\(H=\operatorname{tr}_g h\)、\(\varphi=\dot f\) 及 \(w=H/2-\varphi\)。则
	\begin{equation}\label{g3:eq:perelman-F-variation}
		\frac{\mathrm d}{\mathrm d\varepsilon}\mathcal F
		=\int_M\left[-\langle h,\operatorname{Ric}+\operatorname{Hess}f\rangle
		+w\bigl(s-2\Delta f-|\mathrm df|^2\bigr)\right]e^{-f}\mu_g.
	\end{equation}
\end{lemma}
\begin{proof}
	体积变分给出 \(\partial_\varepsilon(e^{-f}\mu_g)=we^{-f}\mu_g\)，而
	\[
		\partial_\varepsilon|\mathrm df|^2
		=-\langle h,\mathrm df\otimes\mathrm df\rangle
		+2\langle\mathrm df,\mathrm d\varphi\rangle.
	\]
	代入\cref{g3:thm:ricci-scalar-linearization}中的标量曲率变分公式。为整理二阶导数，注意
	\[
		\begin{aligned}
		\Delta(e^{-f})&=-e^{-f}(\Delta f+|\mathrm df|^2),\\
		\operatorname{Hess}(e^{-f})
		&=e^{-f}(-\operatorname{Hess}f+\mathrm df\otimes\mathrm df).
		\end{aligned}
	\]
	闭流形上的分部积分因此给出
	\[
		\begin{aligned}
		\int_M(\Delta H)e^{-f}\mu_g
		&=-\int_MH(\Delta f+|\mathrm df|^2)e^{-f}\mu_g,\\
		\int_M(\delta\delta h)e^{-f}\mu_g
		&=\int_M\langle h,-\operatorname{Hess}f+\mathrm df\otimes\mathrm df\rangle e^{-f}\mu_g,\\
		2\int_M\langle\mathrm df,\mathrm d\varphi\rangle e^{-f}\mu_g
		&=2\int_M\varphi(\Delta f+|\mathrm df|^2)e^{-f}\mu_g.
		\end{aligned}
	\]
	含 \(h(\nabla f,\nabla f)\) 的两项相消。余下含 \(H,\varphi\) 的部分用 \(2\varphi-H=-2w\) 合并，再加密度变分项 \(w(s+|\mathrm df|^2)\)，便得到所列公式。
\end{proof}

为沿 Ricci 流使用这个公式，令正函数 \(u=e^{-f}\) 满足\term[gonge re fangcheng]{共轭热方程}[conjugate heat equation]
\begin{equation}\label{g3:eq:ricci-conjugate-heat}
	\partial_tu=\Delta_{g(t)}u+s_{g(t)}u,
	\qquad\text{即}\qquad
	\partial_tf=\Delta f+|\mathrm df|^2-s.
\end{equation}
它沿 \(t\) 方向是向后的抛物方程。具体的求解方式是在一个光滑流区间 \([0,b]\) 的末端给定光滑正函数 \(u_b\)，再倒转时间：若 \(v(\sigma)=u(b-\sigma)\)，则
\[
	\partial_\sigma v+\Delta_{g(b-\sigma)}v
	=-s_{g(b-\sigma)}v,\qquad v(0)=u_b.
\]
这是正向线性抛物方程。相对于一个固定背景度量，其主系数一致正定，所有系数光滑，故\cref{g3:thm:linear-parabolic-isomorphism,g3:thm:linear-parabolic-smoothing}给出整个闭时间区间上的光滑解。若 \(m=\min_Mu_b>0\)、\(C=\max_{M\times[0,b]}|s|\)，比较函数 \(me^{-C\sigma}\) 与\cref{g3:thm:parabolic-maximum}给出 \(v\ge me^{-C\sigma}>0\)。因此 \(f=-\log u\) 始终有定义。

\ProseParagraphBreak
Ricci 流的体积密度满足 \(\partial_t\mu_g=-s\mu_g\)，所以
\begin{equation}\label{g3:eq:ricci-conjugate-mass}
	\frac{\mathrm d}{\mathrm dt}\int_Mu\mu_g
	=\int_M(\partial_tu-su)\mu_g
	=\int_M\Delta u\,\mu_g=0.
\end{equation}
末端质量为 \(1\) 时，所有较早时刻的质量也为 \(1\)。这样选取权重后，度量的体积变化正好被共轭方程中的零阶项补偿。

\begin{theorem}\label{g3:thm:perelman-F-monotonicity}
	设 \(M\) 为闭光滑流形，\(g(t)\) 为光滑 Ricci 流，实函数 \(f\) 光滑且满足 \(f_t=\Delta_gf+|\mathrm df|_g^2-s_g\)，其中 \(\Delta_g=-\operatorname{div}_g\nabla\)。记 \(\mathcal F(g,f)=\int_M(s_g+|\mathrm df|_g^2)e^{-f}\,\mathrm d\mu_g\)。则
	\begin{equation}\label{g3:eq:perelman-F-monotonicity}
		\frac{\mathrm d}{\mathrm dt}\mathcal F(g(t),f(t))
		=2\int_M|\operatorname{Ric}+\operatorname{Hess}f|^2e^{-f}\mu_g\ge0.
	\end{equation}
	在某一时刻等号成立，当且仅当该时刻 \(\operatorname{Ric}+\operatorname{Hess}f=0\)。
\end{theorem}
\begin{proof}
	令 \(\psi_t\) 为随时间变化的向量场 \(-\nabla_{g(t)}f(t)\) 生成的微分同胚，且令 \(\widetilde g=\psi_t^*g\)、\(\widetilde f=\psi_t^*f\)。闭性保证它在所考虑的紧时间区间上存在。略去拉回记号时，其变分为
	\[
		h=-2(\operatorname{Ric}+\operatorname{Hess}f),
		\qquad \varphi=f_t-|\mathrm df|^2=\Delta f-s.
	\]
	由于 \(\operatorname{tr}\operatorname{Hess}f=-\Delta f\)，有 \(H/2=\Delta f-s=\varphi\)，即 \(w=0\)。将它们代入\cref{g3:lem:perelman-F-variation}，便只剩下所述平方项。泛函及该积分都在同时拉回 \(g,f\) 时不变，故公式也适用于原来的演化。最后，权重处处为正，非负光滑函数的积分为零等价于它处处为零。
\end{proof}

\(\mathcal F\) 的等号对应稳态孤立子。要同时记录收缩尺度，引入正参数 \(\tau\)，定义\term[Perelman W shang]{Perelman \(\mathcal W\) 熵}[Perelman's \(\mathcal W\)-entropy]
\begin{equation}\label{g3:eq:perelman-W-definition}
	\mathcal W(g,f,\tau)\coloneqq
	\int_M\left[\tau(s_g+|\mathrm df|^2)+f-n\right]
	(4\pi\tau)^{-n/2}e^{-f}\mu_g,
\end{equation}
\symidx{\(\mathcal W(g,f,\tau)\)：Perelman 的熵泛函}
其中要求 \(\int_M(4\pi\tau)^{-n/2}e^{-f}\mu_g=1\)。常数缩放给出
\[
	\mathcal W(ag,f,a\tau)=\mathcal W(g,f,\tau),\qquad a>0.
\]
因此 \(\mathcal W\) 与抛物重缩放的尺度相容。

\begin{theorem}\label{g3:thm:perelman-W-monotonicity}
	设 \(M^n\) 闭，\(g_t=-2\operatorname{Ric}\)，\(\tau_t=-1\)、\(\tau>0\)，且
	\[
		f_t=\Delta f+|\mathrm df|^2-s+\frac n{2\tau},
		\qquad u=(4\pi\tau)^{-n/2}e^{-f},\qquad \int_Mu\mu_g=1.
	\]
	记 \(s=\operatorname{Scal}_g\)、\(\Delta=-\operatorname{div}_g\nabla\)，\(\mathcal W(g,f,\tau)=\int_M[\tau(s+|\mathrm df|^2)+f-n]u\,\mathrm d\mu_g\)；所用度量与函数族均联合光滑。则
	\begin{equation}\label{g3:eq:perelman-W-monotonicity}
		\frac{\mathrm d}{\mathrm dt}\mathcal W(g,f,\tau)
		=2\tau\int_M\left|\operatorname{Ric}+\operatorname{Hess}f-\frac1{2\tau}g\right|^2u\mu_g\ge0.
	\end{equation}
\end{theorem}
\begin{proof}
	令 \(\bar f=f+\frac n2\log(4\pi\tau)\)，则 \(u=e^{-\bar f}\)，而 \(\bar f\) 满足式\eqref{g3:eq:ricci-conjugate-heat}。记
	\[
		F(t)=\int_M(s+|\mathrm df|^2)u\mu_g,
		\qquad E(t)=\int_Mfu\mu_g,
		\qquad A=\operatorname{Ric}+\operatorname{Hess}f.
	\]
	因 \(\mathrm d\bar f=\mathrm df\)，上一条定理给出 \(F'=2\int_M|A|^2u\mu_g\)。又由 \(\partial_t(u\mu_g)=(\Delta u)\mu_g\) 及 Laplace 算子的自伴性，
	\[
		\begin{aligned}
		E'&=\int_M(f_t+\Delta f)u\mu_g\\
		&=\int_M\left(2\Delta f+|\mathrm df|^2-s+\frac n{2\tau}\right)u\mu_g
		=-F+\frac n{2\tau}.
		\end{aligned}
	\]
	最后一步使用 \(\int_M(\Delta f)u\mu_g=-\int_M|\mathrm df|^2u\mu_g\)。于是 \(\mathcal W=\tau F+E-n\) 满足
	\[
		\mathcal W'=2\tau\int_M|A|^2u\mu_g-2F+\frac n{2\tau}.
	\]
	而 \(\int_M\operatorname{tr}A\,u\mu_g=\int_M(s-\Delta f)u\mu_g=F\)。展开 \(2\tau|A-g/(2\tau)|^2\) 并积分，恰得到上一行。
\end{proof}

因此 \(\mathcal W\) 导数为零恰好对应 \(\operatorname{Ric}+\operatorname{Hess}f=g/(2\tau)\)，即收缩孤立子方程。上面的证明也说明，使用熵公式时须把 \(f\) 与共轭热方程一起演化；任意固定一个权函数并不能得到同样的单调性。

\subsection{有限时间非塌缩}
\label{g3:subsec:2704:04}
曲率界控制局部弯曲，体积下界控制一个球是否在较少的方向上变得极薄。先固定本节使用的时空版本。对 \(r^2\le t<T\)，定义\term[paowu qiu]{抛物球}[parabolic ball]
\begin{equation}\label{g3:eq:ricci-parabolic-ball}
	P(x,t,r)\coloneqq B_{g(t)}(x,r)\times[t-r^2,t].
\end{equation}
这里整个时间区间都使用最终时刻的同一个点集 \(B_{g(t)}(x,r)\)，而在该点集上计算曲率范数时使用各时刻自己的度量。

\begin{definition}\label{g3:def:kappa-noncollapsed}
	设 \(g(t)\)、\(0\le t<T\)，为光滑无边界 \(n\) 维流形上的 Ricci 流，记 \(P(x,t,r)=B_{g(t)}(x,r)\times[t-r^2,t]\)。
	给定 \(\kappa>0\)、\(\rho>0\)。若对每个 \(0<r<\rho\) 及 \(r^2\le t<T\)，只要
	\[
		|\operatorname{Rm}_{g(a)}|(y)\le r^{-2}
		\quad\text{对所有 }(y,a)\in P(x,t,r)
	\]
	就有
	\begin{equation}\label{g3:eq:ricci-kappa-volume}
		\operatorname{Vol}_{g(t)}B_{g(t)}(x,r)\ge\kappa r^n,
	\end{equation}
	则称该 Ricci 流在小于 \(\rho\) 的尺度上\term[kappa feitasuo]{\(\kappa\) 非塌缩}[\(\kappa\)-noncollapsed]。
\end{definition}

\begin{theorem}[有限时间非塌缩定理]\label{g3:thm:ricci-no-local-collapsing}
	设 \(M^n\) 为闭光滑流形，\(g(t)\) 为 \([0,T)\) 上的光滑 Ricci 流，\(0<T<\infty\)。对任意固定 \(\rho>0\)，存在仅依赖 \(g(0),T,\rho\) 的 \(\kappa>0\)，使该流在小于 \(\rho\) 的尺度上 \(\kappa\) 非塌缩。
\end{theorem}

这是 Perelman 的\term[jubu feitasuo dingli]{局部非塌缩定理}[no local collapsing theorem]在闭流形上的版本，作为本节的外部定理使用；精确的时空定义和定理见 \cite[Definition~26.1、Theorem~26.2，第2645页]{KleinerLott2008Perelman}。该文的定理还适用于初始完备、曲率有界且单射半径有正下界的非紧流形，并要求每个较短闭时间区间上的曲率一致有界。闭流形和光滑流自动满足这些输入条件。因此这里不要求初始曲率具有正性，也不要求终止时刻前的全局曲率有统一上界。

\ProseParagraphBreak
非塌缩常数在抛物重缩放下保持不变。事实上，将 \(g\) 乘以 \(Q\)、时间乘以 \(Q\) 后，半径 \(r\) 变成 \(Q^{1/2}r\)，曲率界变成 \((Q^{1/2}r)^{-2}\)，且
\[
	\frac{\operatorname{Vol}_{Qg}B_{Qg}(x,Q^{1/2}r)}{(Q^{1/2}r)^n}
	=\frac{\operatorname{Vol}_gB_g(x,r)}{r^n}.
\]
应用于式\eqref{g3:eq:ricci-backward-rescaled-bound}，取 \(r_j=Q_j^{-1/2}\)。当 \(j\) 充分大时，\(r_j<\rho\)、\(r_j^2\le t_j\)，且相关整个抛物球上的曲率不超过 \(Q_j=r_j^{-2}\)。于是
\begin{equation}\label{g3:eq:rescaled-unit-ball-volume}
	\operatorname{Vol}_{g_j(0)}B_{g_j(0)}(x_j,1)\ge\kappa.
\end{equation}
这正是奇点放大时需要的尺度一致体积信息。与导数估计和局部几何紧性结合时，它能排除因体积退化而丢失维数的情形。

\ProseParagraphBreak
曲率界本身不能代替这个输入。例如让一个平坦 \(n\) 维环面中的一个圆因子长度为 \(\varepsilon\)，其余因子长度固定且大于 \(4\)。曲率恒为零，单位球的体积却至多为 \(\varepsilon\) 乘以一个固定的 \((n-1)\) 维球体积，因而随 \(\varepsilon\downarrow0\) 趋于零。不同的初始度量不能共用一个与初始几何无关的 \(\kappa\)。同样，定理保留了有限的时间上界 \(T\) 和尺度上界 \(\rho\)，并未给出在所有时间、所有尺度上统一有效的体积下界。
% 纲要条目排为编号小节标题，并进入目录与书签；说明段仍保留于纲要注释中。

% !TeX root = ../../Book3.tex
% 纲要标题：### 27.5 曲面与齐性算例
% 纲要位置：几何学计划纲要.md，第 3272 行。
% 正文状态：初稿；尚未正式审核。
% 后续写作任务：本节只处理所列曲面与齐性流模型，并准确声明熵、非塌缩、三维手术与几何化所用的外部结果；拓扑分解和几何化陈述与第二卷附录A第4节分工，双曲商与 Mostow–Prasad 刚性另见本卷附录D的规划，不把尚未写成的目标当作本节已证输入。

% 纲要内容开始
% 
% * 常曲率曲面的显式演化
% * 归一化和非归一化流
% * 齐性度量的常微分方程
% * 几何化和手术理论的明确边界
% 
% ---
% 
% 纲要内容结束

\section{曲面与齐性算例}
\label{g3:sec:2705}

方程的一般理论说明何时能求解、何时可以延拓；对称模型则把这些条件转成可以计算的长度与曲率。本节先求曲面上的尺度变化，再用体积归一化比较不同时间的几何。齐性条件把空间偏微分方程约化为有限维常微分方程，但保留了各向异性：不同方向可能同时增长和缩短。最后说明从这些光滑模型走向三维手术与几何化还需要哪些整体结果。

\subsection{常曲率曲面的显式演化}
\label{g3:subsec:2705:01}
在二维，Ricci 张量由 Gauss 曲率完全确定：\(\operatorname{Ric}=Kg\)、\(s=2K\)。因此 Ricci 流保持初始度量的点态共形类。这个简化来自维数二，而不是一般 Ricci 流的性质。

\begin{proposition}\label{g3:prop:ricci-flow-surface-conformal}
	设 \(M^2\) 为闭光滑曲面，\(g_0\) 为光滑度量，\(g(t)\) 是以 \(g_0\) 为初值的光滑 Ricci 流。则存在实值光滑函数 \(u(x,t)\)，使 \(g(t)=e^{2u(x,t)}g_0\)、\(u(x,0)=0\)，并且
	\begin{equation}\label{g3:eq:ricci-flow-surface-conformal}
		u_t=-e^{-2u}(K_0+\Delta_{g_0}u).
	\end{equation}
	若 \(K_0=k\) 为常数，则在 \(1-2kt>0\) 的区间上
	\begin{equation}\label{g3:eq:ricci-flow-constant-surface}
		g(t)=(1-2kt)g_0,\qquad K(t)=\frac{k}{1-2kt}.
	\end{equation}
\end{proposition}
\begin{proof}
	对固定 \(x\)，方程 \(g_t=-2K(x,t)g\) 是矩阵各分量乘同一标量的常微分方程，故 \(g=e^{2u}g_0\)，其中 \(u(x,t)=-\int_0^tK(x,a)\,\mathrm da\)。\Cref{g3:cor:conformal-scalar-power}在维数二给
	\(K_{e^{2u}g_0}=e^{-2u}(K_0+\Delta_{g_0}u)\)，代回 \(u_t=-K\) 即得第一式。常曲率时，\(u(t)=\tfrac12\log(1-2kt)\) 满足该方程；也可直接由\cref{g3:exm:ricci-flow-einstein-scaling}验证。\Cref{g3:thm:ricci-flow-short-time}的唯一性使它等于给定初值的流。
\end{proof}

当 \(k>0\) 时，面积线性减小，曲率在 \(t=1/(2k)\) 发散；当 \(k=0\) 时度量静止；当 \(k<0\) 时度量对所有正时间膨胀，而曲率从负侧趋于零。后一个极限只说明未归一化曲率变小，长度尺度却趋于无穷，因而不能直接把度量极限称为一个固定的平坦度量。

\begin{proposition}\label{g3:prop:ricci-flow-surface-area}
	设 \(M\) 为闭定向光滑曲面，\(g(t)\) 是其光滑 Ricci 流，\(A(t)=\operatorname{Vol}_{g(t)}(M)\)。则在光滑存在期间
	\[
		A(t)=A(0)-4\pi\chi(M)t.
	\]
	特别地，若 \(\chi(M)>0\)，光滑流的最大正存在时间 \(T_{\max}\) 满足 \(T_{\max}\le A(0)/(4\pi\chi(M))\)。
\end{proposition}
\begin{proof}
	由体积变分 \(\partial_t\mu_g=-s\mu_g=-2K\mu_g\)，有 \(A'=-2\int_MK\mu_g\)。\Cref{g3:thm:gauss-bonnet-chern}的二维情形给 \(\int_MK\mu_g=2\pi\chi(M)\)，所以导数为上述常数。积分得到结论。若流在给定上界之后仍光滑，公式将产生非正面积，与正定度量矛盾。
\end{proof}

面积公式适用于非恒定初始曲率，但它本身没有证明在这个上界之前不会出现曲率奇点。一般曲面的长期归一化行为还需要针对其共形方程的估计。

\subsection{归一化和非归一化流}
\label{g3:subsec:2705:02}
整体缩放可以遮住形状变化。例如正曲率球面的体积趋于零，但除去尺度后它的形状始终不变。固定体积提供了比较形状的一种统一方式。

\begin{definition}\label{g3:def:normalized-ricci-flow}
	设 \(M^n\) 闭，\(n\ge2\)，置
	\[
		r(g)=\frac{1}{\operatorname{Vol}_g(M)}\int_Ms_g\,\mu_g.
	\]
	度量族若满足
	\begin{equation}\label{g3:eq:normalized-ricci-flow}
		\partial_\tau g=-2\operatorname{Ric}_g+\frac2n r(g)g,
	\end{equation}
	则称为保持体积的\term[guiyihua Ricci liu]{归一化 Ricci 流}[normalized Ricci flow]。
\end{definition}

\begin{proposition}\label{g3:prop:ricci-flow-normalization}
	式\eqref{g3:eq:normalized-ricci-flow}保持总体积。若 \(g(t)\) 是闭 \(n\) 维流形上的光滑未归一化 Ricci 流，记 \(V(t)=\operatorname{Vol}_{g(t)}(M)\)，令
	\begin{equation}\label{g3:eq:ricci-flow-normalizing-change}
		a(t)=\left(\frac{V(0)}{V(t)}\right)^{2/n},\qquad
		\tau(t)=\int_0^t a(s)\,\mathrm ds,\qquad
		\widehat g(\tau(t))=a(t)g(t).
	\end{equation}
	则 \(\widehat g\) 满足归一化流，体积恒为 \(V(0)\)。闭流形上的 Einstein 度量是归一化方程的静止解。
\end{proposition}
\begin{proof}
	取归一化方程的一半迹，得到 \(\partial_\tau\mu_g=(-s+r)\mu_g\)，积分为零。对未归一化流，\(V'=-r(g)V\)，所以 \(a'/a=2r(g)/n\)。因 \(\tau'=a>0\)，可把 \(t\) 作为 \(\tau\) 的函数，并计算
	\[
		\partial_\tau\widehat g
		=\frac{a'}{a^2}\widehat g-2\operatorname{Ric}_{g(t)}
		=\frac2n r(\widehat g)\widehat g-2\operatorname{Ric}_{\widehat g}.
	\]
	第二步用了常数缩放下 \(r(ag)=a^{-1}r(g)\) 以及 Ricci 协变二张量不变。体积为 \(a^{n/2}V=V(0)\)。若 \(\operatorname{Ric}=\lambda g\)，则 \(r=n\lambda\)，方程右侧为零。
\end{proof}

\begin{example}\label{g3:exm:ricci-normalized-time}
	对 \(g(t)=(1-2\lambda t)g_0\)，有 \(a(t)=(1-2\lambda t)^{-1}\)，故 \(\widehat g=g_0\)。当 \(\lambda\ne0\) 时
	\[
		\tau(t)=-\frac{1}{2\lambda}\log(1-2\lambda t).
	\]
	正 \(\lambda\) 的有限消失时间被变为 \(\tau\to\infty\)；负 \(\lambda\) 的 \(t\to\infty\) 也对应 \(\tau\to\infty\)。所以“有限时奇点”必须说明使用哪一种时间和归一化。
\end{example}

在曲面上，固定面积 \(A_0\) 时 \(r=4\pi\chi(M)/A_0\) 为常数。若 \(g=e^{2u}g_0\)，归一化方程相应变成
\[
	 u_\tau=-e^{-2u}(K_0+\Delta_{g_0}u)+\frac{2\pi\chi(M)}{A_0}.
\]
常曲率度量是这个方程的平衡点。是否由任意初始度量收敛到平衡点，是一个另需先验估计与紧性的动力学问题。

\subsection{齐性度量的常微分方程}
\label{g3:subsec:2705:03}
若一组固定的微分同胚在初时刻保持度量，自然性与唯一性会使这组对称性沿流保存。对传递群作用，整个度量便由一个点处的有限个数决定。

\begin{proposition}\label{g3:prop:ricci-flow-homogeneous-reduction}
	设闭光滑流形 \(M\) 上有 Lie 群 \(G\) 的光滑传递作用，\(g_0\) 是 \(G\) 不变度量。则其光滑 Ricci 流始终 \(G\) 不变。固定 \(p\in M\)，令 \(H\) 为其稳定子，\(\mathcal P\) 为 \(\mathrm T_pM\) 上 \(H\) 不变正定内积组成的开锥。流方程在这个有限维锥上化为
	\[
		b'(t)=-2\operatorname{Ric}_{g_{b(t)}}|_p,
	\]
	其中 \(g_b\) 是由 \(b\) 沿群作用延拓的唯一 \(G\) 不变度量，右端是 \(b\) 的光滑函数。
\end{proposition}
\begin{proof}
	对每个固定群元素 \(a\)，\(a^*g(t)\) 与 \(g(t)\) 解同一个初值问题，由\cref{g3:thm:ricci-flow-short-time}相同，故对称性保持。\(H\) 不变性恰保证把 \(b\) 从 \(p\) 送到 \(ap\) 时不依赖代表 \(a\)；这构成上述一一对应。
	取有限个固定局部截面来表达群作用，所得 \(g_b\) 的坐标分量对 \(b\) 线性，空间导数也光滑依赖 \(b\)。联络与曲率由分量、其导数和逆矩阵构成，因此 Ricci 对 \(b\) 光滑。其自然性保证它仍是 \(G\) 不变张量。于是正定性保留期间，有限维方程与原 Ricci 方程等价。
\end{proof}

下面直接在三维幂零 Lie 群上计算。它展示各向异性的增长，而不只是在 Einstein 度量上整体缩放。取 \(\mathbb R^3\) 上的左不变余标架
\[
	\alpha=\mathrm dx,\quad\beta=\mathrm dy,\quad
	\gamma=\mathrm dz-x\,\mathrm dy,
	\qquad g=A\alpha^2+B\beta^2+C\gamma^2,
\]
其中 \(A,B,C>0\) 在空间中为常数。对应群乘法为
\((x,y,z)(x',y',z')=(x+x',y+y',z+z'+xy')\)。对偶场
\(X=\partial_x\)、\(Y=\partial_y+x\partial_z\)、\(Z=\partial_z\) 满足 \([X,Y]=Z\)，其余括号为零。

\begin{proposition}\label{g3:prop:ricci-flow-nil-model}
	上述度量族满足 Ricci 流，当且仅当
	\begin{equation}\label{g3:eq:ricci-flow-nil-ode}
		A'=\frac CB,\qquad B'=\frac CA,\qquad
		C'=-\frac{C^2}{AB}.
	\end{equation}
	对任意正初值，置 \(q_0=C_0/(A_0B_0)\)、\(L(t)=1+3q_0t\)，其解为
	\begin{equation}\label{g3:eq:ricci-flow-nil-solution}
		A=A_0L^{1/3},\qquad B=B_0L^{1/3},\qquad
		C=C_0L^{-1/3},\qquad t\ge0.
	\end{equation}
\end{proposition}
\begin{proof}
	取正交标架 \(e_1=X/\sqrt A\)、\(e_2=Y/\sqrt B\)、\(e_3=Z/\sqrt C\)，令 \(c=\sqrt{C/(AB)}\)，则 \([e_1,e_2]=ce_3\)。左不变场的 Koszul 公式为
	\[
		2\langle\nabla_{e_i}e_j,e_k\rangle
		=\langle[e_i,e_j],e_k\rangle
		-\langle[e_j,e_k],e_i\rangle+\langle[e_k,e_i],e_j\rangle.
	\]
	故非零联络系数为
	\[
	\begin{aligned}
		\nabla_1e_2&=\tfrac c2e_3,&\nabla_2e_1&=-\tfrac c2e_3,\\
		\nabla_1e_3=\nabla_3e_1&=-\tfrac c2e_2,&
		\nabla_2e_3=\nabla_3e_2&=\tfrac c2e_1.
	\end{aligned}
	\]
	各个 \(\nabla_ie_i\) 为零。代入曲率定义得到
	\[
		K_{12}=-\frac{3c^2}{4},\quad K_{13}=K_{23}=\frac{c^2}{4},
		\qquad (\operatorname{Ric}_{ij})=
		\operatorname{diag}(-c^2/2,-c^2/2,c^2/2).
	\]
	例如 \(R(e_1,e_2)e_2=-\nabla_2(c e_3/2)-c\nabla_3e_2=-3c^2e_1/4\)；另两平面直接给 \(c^2/4\)。由所列联络也可核得混合 Ricci 分量为零。回到固定的 \(X,Y,Z\) 标架，三个分量分别乘 \(A,B,C\)，代入 \(g_t=-2\operatorname{Ric}\) 即得方程。
	令 \(q=C/(AB)\)，则 \((\log A)'=(\log B)'=q\)、\((\log C)'=-q\)，从而 \(q'=-3q^2\)。因此 \(q=q_0/L\)，分别积分得到式\eqref{g3:eq:ricci-flow-nil-solution}。
\end{proof}

整数坐标组成此群的离散子群 \(\Gamma\)。先用整数左平移把 \(x,y\) 调到 \([0,1]\)，再用中心的整数平移把 \(z\) 调到 \([0,1]\)，可见 \(\Gamma\backslash\mathbb R^3\) 紧。左不变度量对 \(\Gamma\) 不变，所算流因此下降到这个紧商。商在局部保留齐性公式；不必要求左平移群整体作用在左陪集商上。

\ProseParagraphBreak
在该紧商上，两水平长度按 \(L^{1/6}\) 增长，中心圆纤维的长度按 \(L^{-1/6}\) 缩短；体积按 \(L^{1/6}\) 增长，曲率按 \(L^{-1}\) 衰减。它对全部正时间光滑存在，却始终不是各方向同一比例的膨胀。这也说明有限时间、有限尺度的非塌缩结论不会自动提供跨越所有大时间的统一长度下界。

\subsection{几何化和手术理论的明确边界}
\label{g3:subsec:2705:04}
光滑 Ricci 流在固定流形上改变度量。当曲率在有限时刻发散时，\cref{g3:thm:ricci-flow-curvature-continuation}说明继续光滑求解的障碍已经出现。三维手术方法在受控的高曲率颈部切开，并粘入标准帽，随后在新的时间片上重新开始 Ricci 流。这里的颈部拓扑为 \(S^2\times(-1,1)\)，切割与补帽沿二球面进行；第二卷附录 A.3 的 Dehn 填充沿环面边界进行，是另一种拓扑操作。

\ProseParagraphBreak
\term[dai shoushu de Ricci liu]{带手术的 Ricci 流}[Ricci flow with surgery]由一系列光滑 Ricci 流区间、奇异时刻的保留区域、两侧保留区域之间的等距识别，以及新添或移除的部分组成。可精确对照文献\cite[定义68.1，第2746--2747页；第73节]{KleinerLott2008Perelman}。为执行三维几何化中的手术，任意这样的拼接还不够：须规定颈部在缩放后接近标准圆柱的程度、切断尺度、标准帽的几何，并验证手术后曲率夹挤与非塌缩估计仍然有效。

\begin{theorem}\label{g3:thm:ricci-surgery-existence-external}
	设 \(M^3\) 是闭、连通、可定向光滑三维流形，\(g_0\) 为光滑黎曼度量。适当常数缩放初值后，可以选择手术参数，按文献的 \((r,\delta)\) 截断构造得到定义于全部正时间的带手术 Ricci 流。每个光滑时间区间上满足 \(g_t=-2\operatorname{Ric}\)，手术按规定的圆柱颈部切开并补帽，允许移去规定类型的连通分支；若所有分支已经消失，此后时间片取为空集。
\end{theorem}

该存在性取自\cite[定义77.1、命题77.2，第2770--2771页；证明见第78--80节]{KleinerLott2008Perelman}。其中标准化初值要求 \(|\mathrm{Rm}|\le1\)，且每个单位球的体积至少为 Euclid 单位球体积的一半。闭光滑初值可以通过放大度量达到这两个条件：曲率范数按缩放因子的逆数变化，而紧性使小测地球的体积比一致趋于 Euclid 值。命题给出随时间变化的颈部与手术精度参数，使构造能够继续；它是本节采用的外部深层结果。

\ProseParagraphBreak
从手术存在到几何化，还要控制无穷时间的厚部与薄部：厚部的适当缩放极限具有双曲几何，薄部的塌缩须转化为图流形的拓扑结构，并证明用于分解的环面具有所需的不可压缩性。这些步骤分别见所引文献第89--92节。第二卷附录 A.4.4 已准确陈述并外引紧可定向三维流形的几何化及球面化版本；本节说明该结论中 Ricci 流所承担的分析工作，不重新把这些整体结论作为本章已证定理。

\begin{table}[htbp]
	\centering
	\small
	\caption{本章几种延续与比较的对象}
	\label{g3:tab:ricci-flow-scope}
	\begin{tabularx}{\textwidth}{@{}l X X@{}}
		\toprule
		情形 & 保留的对象 & 所得结论 \\
		\midrule
		闭光滑初值 & 固定流形上的度量 & 短时存在、唯一性；有限时曲率有界则光滑延拓 \\
		体积归一化 & 同一度量流的尺度与时间变换 & 保持总体积；终止时间随重参数改变 \\
		三维规定手术 & 光滑区间及手术处的保留区域 & 可选择参数作全时构造，允许时间片消失 \\
		几何化分解 & 三维流形的拓扑块及其模型几何 & 还使用长期厚薄分析、塌缩及不可压缩性 \\
		\bottomrule
	\end{tabularx}
\end{table}

\cref{g3:tab:ricci-flow-scope}中的光滑唯一性固定初始度量与参数化流形；手术构造还包含切断与补帽的选择。本章所引手术存在定理不提供这些选择之间的唯一性，也不由球面或齐性模型的计算推出一般三维流形的长期行为。
% 纲要条目排为编号小节标题，并进入目录与书签；说明段仍保留于纲要注释中。


\begin{chaptersummary}
Ricci 流在微分同胚方向上退化。固定背景度量后，DeTurck 向量场使方程的最高阶部分成为 \(g^{ab}\bar\nabla_a\bar\nabla_b g\)。强抛物短时解经 \(-W\) 的流拉回，恢复 Ricci 流；反向的调和映射热方程把任意 Ricci 流送入同一规范，给出固定初始度量的唯一性。

\ProseParagraphBreak
完整曲率在固定正交丛上满足 \((\partial_t+\widehat\Delta)\mathcal R=2(\mathcal R^2+\mathcal R^\#)\)。非负曲率算子的保持来自反应项对非负锥的相容性；标量最大值原理则控制标量曲率及曲率导数的长度。逐阶使用负梯度项，可以把曲率界提升为所有正时间的导数界，进而得到光滑终点度量与有限时延拓。

\ProseParagraphBreak
抛物重缩放同时改变长度和时间单位，孤立子描述在缩放与微分同胚作用下形状不变的模型。共轭热权重使 \(\mathcal F\)、\(\mathcal W\) 的导数成为平方积分，等号分别对应稳态和收缩孤立子。有限时间非塌缩定理提供曲率尺度上的体积下界；取得光滑放大极限还要配合导数估计与几何紧性。

\ProseParagraphBreak
曲面的共形方程、Einstein 度量的缩放和 Nil 度量的各向异性演化说明，曲率、长度与体积可以有不同的变化趋势。体积归一化只是尺度和时间的变换；三维手术则改变时间片，并需受控颈部、标准帽和长期分析。其存在性与几何化的整体结论按末节所列来源使用。
\end{chaptersummary}

\BookExercises

\begin{exercise}\label{g3:ex:ricci-einstein-products}
设 \((M^p,g_1)\)、\((N^q,g_2)\) 为闭 Einstein 流形，\(\operatorname{Ric}_{g_i}=\lambda_i g_i\)，且 \(a_0,b_0>0\)。求从 \(a_0g_1\oplus b_0g_2\) 出发的 Ricci 流、标量曲率及体积。判断何时该流仅作整体缩放。
\end{exercise}
\begin{solution}
乘积联络不混合两个因子，常数缩放也不改变各自的 Ricci 协变二张量。因此
\[
g(t)=a(t)g_1\oplus b(t)g_2,\qquad
a(t)=a_0-2\lambda_1t,\quad b(t)=b_0-2\lambda_2t.
\]
解在两个系数都为正的最大正时间区间上存在。若至少一个 \(\lambda_i>0\)，终点是相应正数 \(a_0/(2\lambda_1)\)、\(b_0/(2\lambda_2)\) 中的最小者；若均非正，终点为无穷。标量曲率与体积为
\[
s=\frac{p\lambda_1}{a}+\frac{q\lambda_2}{b},\qquad
V=a^{p/2}b^{q/2}V_{g_1}(M)V_{g_2}(N).
\]
有限终点处某个正 Einstein 因子的 Ricci 范数发散，故确是光滑最大时间。比例 \(a/b\) 恒定当且仅当 \(\lambda_1/a_0=\lambda_2/b_0\)；此时整个初始度量为 Einstein，得到整体缩放，否则两个因子的相对尺度发生变化。
\end{solution}

\begin{exercise}\label{g3:ex:ricci-time-reparametrization}
在平坦环面 \((\mathbb R/2\pi\mathbb Z)^n\) 上取 \(g_0=\sum_i\mathrm dx_i^2\)。令 \(\psi_t(x)=(x_1+t\sin x_1,x_2,\ldots,x_n)\)，其中 \(|t|<1\)。证明 \(\psi_t^*g_0\) 虽然每一时刻平坦，却不是 Ricci 流。对一般静止 Ricci 平坦度量，写出时间依赖拉回仍为 Ricci 流的条件。
\end{exercise}
\begin{solution}
因为 \(1+t\cos x_1>0\)，\(\psi_t\) 是环面的微分同胚，而
\[
\psi_t^*g_0=(1+t\cos x_1)^2\mathrm dx_1^2+\sum_{i=2}^n\mathrm dx_i^2.
\]
它的 Ricci 张量为零，但时间导数在 \(t=0\) 为 \(2\cos x_1\,\mathrm dx_1^2\ne0\)，故不满足流方程。一般地，若 \(X_t=\dot\psi_t\circ\psi_t^{-1}\)，则 \(\partial_t(\psi_t^*g_0)=\psi_t^*\mathcal L_{X_t}g_0\)。由于右边的 Ricci 张量仍为零，该族为 Ricci 流当且仅当每个 \(X_t\) 都是 \(g_0\) 的 Killing 场。
\end{solution}

\begin{exercise}\label{g3:ex:ricci-flat-linearization}
在 \(T^n=(\mathbb R/2\pi\mathbb Z)^n\)、\(n\ge2\)，上取 \(g_0=\sum_i\mathrm dx_i^2\)。在这一平坦环面上，以 \(g_0\) 为背景线性化 Ricci--DeTurck 方程。求解初值为光滑对称二张量 \(h_0\) 的线性方程，并与未加规范项时的微分同胚方向比较。
\end{exercise}
\begin{solution}
平坦背景下曲率零，Lichnerowicz Laplace 算子逐分量等于 \(-\sum_a\partial_a^2\)，所以线性方程为 \(h_t=\sum_a\partial_a^2h\)。写 Fourier 展开
\[
h_0(x)=\sum_{k\in\mathbb Z^n}H_k e^{\mathrm i k\cdot x},\qquad
h(x,t)=\sum_{k\in\mathbb Z^n}e^{-|k|^2t}H_k e^{\mathrm i k\cdot x}.
\]
初值光滑使系数迅速衰减，故逐项微分成立。零频率不变，非零频率衰减。另一方面，对任意周期向量场 \(X\)，Ricci 张量的自然性给 \(D\operatorname{Ric}_{g_0}(\mathcal L_Xg_0)=\mathcal L_X\operatorname{Ric}_{g_0}=0\)。未加规范时这些方向的线性演化静止；DeTurck 项使其中非零频率也服从热方程。
\end{solution}

\begin{exercise}\label{g3:ex:ricci-average-scalar}
对闭 \(n\) 维流形的体积归一化 Ricci 流，记 \(V=\operatorname{Vol}_{g(t)}M\)、\(s=\operatorname{Scal}_{g(t)}\)，\(\langle v\rangle=V^{-1}\int_Mv\mu_g\)、\(r=\langle s\rangle\)，并令 \(\operatorname{Ric}^{\circ}=\operatorname{Ric}-(s/n)g\)。证明
\[
r'=2\langle|\operatorname{Ric}^{\circ}|^2\rangle
-\frac{n-2}{n}\langle(s-r)^2\rangle.
\]
解释为何二维的平均标量曲率恒定，而高维时仅凭此式不能断言它非减。
\end{exercise}
\begin{solution}
归一化项是空间常数倍度量，其对标量曲率的变分为 \(-(2/n)rs\)，因此
\[
s_\tau=-\Delta s+2|\operatorname{Ric}|^2-\frac2nrs,
\qquad \partial_\tau\mu_g=(-s+r)\mu_g.
\]
体积不变且 \(\int\Delta s=0\)，故
\[
r'=2\langle|\operatorname{Ric}|^2\rangle-\langle s^2\rangle
+\left(1-\frac2n\right)r^2.
\]
代入 \(|\operatorname{Ric}|^2=|\operatorname{Ric}^{\circ}|^2+s^2/n\) 并用 \(\langle(s-r)^2\rangle=\langle s^2\rangle-r^2\)，得到所求式。在二维 \(\operatorname{Ric}^{\circ}=0\)，两个项都为零，也与 Gauss--Bonnet 和固定面积相符。高维时两个非负量带相反符号，若无进一步关系，这个恒等式不提供单一符号。
\end{solution}

\begin{exercise}\label{g3:ex:ricci-curvature-rate}
设闭 Ricci 流的最大正存在时间 \(T\) 有限。证明存在只依赖维数的 \(c_n>0\)，使
\[
\max_M|\mathrm{Rm}|(t)\ge\frac{c_n}{T-t},\qquad 0\le t<T.
\]
提示：对 \(F=|\mathrm{Rm}|^2\) 使用 \((\partial_t+\Delta)F\le C_nF^{3/2}\)，再使用延拓判据。
\end{exercise}
\begin{solution}
固定 \(t_0<T\)，记 \(K_0=\max|\mathrm{Rm}|(t_0)\)。若 \(K_0=0\)，平坦静止解与唯一性使流永远存在，矛盾。设 \(K_0>0\)，常微分方程 \(b'=C_nb^{3/2}\)、\(b(t_0)=K_0^2\) 给
\[
b(t)=\frac{K_0^2}{\bigl(1-(C_n/2)K_0(t-t_0)\bigr)^2}.
\]
在分母为正的共同时间内，标量比较给 \(F\le b\)：可对 \(F-b\) 使用局部有界的差商系数，或先把初值提高再取极限。若 \(T-t_0<2/(C_nK_0)\)，则这个上界一直到 \(T\) 都有限；加上 \([0,t_0]\) 上的光滑界，\cref{g3:thm:ricci-flow-curvature-continuation}将流延拓过 \(T\)，矛盾。所以 \(T-t_0\ge2/(C_nK_0)\)，即取 \(c_n=2/C_n\) 得到结论。
\end{solution}

\begin{exercise}\label{g3:ex:ricci-three-eigenvalues}
在三维欧氏空间的 \(\Lambda^2\) 上考虑反应常微分方程 \(\mathcal R'=2(\mathcal R^2+\mathcal R^\#)\)。若初始算子在 \(E_{23},E_{31},E_{12}\) 中对角，求三个对角元 \(a,b,c\) 的方程，并直接证明非负正交卦限保持。说明这个计算与空间依赖曲率的保持定理之间还相差哪一步。
\end{exercise}
\begin{solution}
所取正交基的循环括号具有单位长度。代入 \(\mathcal R^\#\) 的定义，其非对角元为零，对角元为 \(bc,ca,ab\)，所以对角子空间不变，且
\[
a'=2(a^2+bc),\qquad b'=2(b^2+ca),\qquad c'=2(c^2+ab).
\]
在三个变量均非负时，右端均非负。严格正初值的三个量在解存在期间非减，故不可能离开正卦限；非负初值用严格正初值逼近，由常微分方程的连续依赖在任意共同紧时间区间取极限，得到非负性。对于空间依赖的曲率，方程还含联络 Laplace 项，特征向量也随位置变化；正文的固定内积丛和首次零方向论证正用于处理这一步，不能只逐点解上述反应方程。
\end{solution}

\begin{exercise}\label{g3:ex:ricci-compact-steady}
设 \((M,g,f)\) 为闭、连通的梯度稳态 Ricci 孤立子，即 \(\operatorname{Ric}+\operatorname{Hess}f=0\)。证明 \(\operatorname{Ric}=0\)，且 \(f\) 为常数。
\end{exercise}
\begin{solution}
由\cref{g3:prop:soliton-selfsimilar-flow}，\(g(t)=\phi_t^*g\)，其中 \(\dot\phi_t=\nabla f\circ\phi_t\)。标量曲率也随之拉回，故在 \(t=0\) 有 \(s_t=\langle\mathrm ds,\mathrm df\rangle\)。与标量演化比较得
\[
\Delta s+\langle\mathrm ds,\mathrm df\rangle=2|\operatorname{Ric}|^2.
\]
对权 \(e^{-f}\mu_g\) 积分时，分部积分给
\(\int(\Delta s)e^{-f}\mu_g=-\int\langle\mathrm ds,\mathrm df\rangle e^{-f}\mu_g\)。因此 \(\int2|\operatorname{Ric}|^2e^{-f}\mu_g=0\)，由权正性知 Ricci 处处为零。孤立子方程于是给 \(\operatorname{Hess}f=0\)，取迹得 \(\Delta f=0\)。最后 \(\int|\mathrm df|^2=\int f\Delta f=0\)，连通性使 \(f\) 为常数。
\end{solution}

\begin{exercise}\label{g3:ex:ricci-gaussian-entropy}
在 Euclid 空间的静止 Ricci 流上，取 \(\tau=T-t>0\)、\(f(x,t)=|x|^2/(4\tau)\) 及 \(u=(4\pi\tau)^{-n/2}e^{-f}\)。直接验证 \(u\) 的共轭热方程及质量归一化，并计算 \(\mathcal W\)。
\end{exercise}
\begin{solution}
有 \(\Delta f=-n/(2\tau)\)、\(|\mathrm df|^2=|x|^2/(4\tau^2)=f/\tau\)，而 \(f_t=f/\tau\)，故 \(f_t=\Delta f+|\mathrm df|^2+n/(2\tau)\)。因此 \(u_t=\Delta u\)，与零标量曲率的共轭热方程一致。逐坐标 Gaussian 积分给 \(\int u\,\mathrm dx=1\) 和 \(\int|x|^2u\,\mathrm dx=2n\tau\)，于是
\[
\mathcal W=\int_{\mathbb R^n}(2f-n)u\,\mathrm dx=0.
\]
同时 \(\operatorname{Hess}f=g_{\mathrm E}/(2\tau)\)，单调公式中的平方项为零。这里在非紧空间直接计算收敛积分，所得结论不需把闭流形定理未经说明地移用到无穷远。
\end{solution}

\begin{exercise}\label{g3:ex:ricci-einstein-F}
设 \((M^n,g_0)\) 为闭 Einstein 流形，\(\operatorname{Ric}_{g_0}=\lambda g_0\)，\(g(t)=a(t)g_0\)、\(a=1-2\lambda t>0\)，并令 \(f(t)=\log\operatorname{Vol}_{g(t)}M\)。证明 \(e^{-f}\) 解共轭热方程且质量为一，并直接核对 \(\mathcal F\) 的单调公式。
\end{exercise}
\begin{solution}
标量曲率为 \(s=n\lambda/a\)，体积为 \(V=a^{n/2}V_0\)，所以 \(f_t=V'/V=-n\lambda/a=-s\)。由于 \(f\) 在空间中为常数，这正是 \(f_t=\Delta f+|\mathrm df|^2-s\)，归一化也由定义成立。于是
\[
\mathcal F=\frac{n\lambda}{a},\qquad
\mathcal F'=\frac{2n\lambda^2}{a^2}.
\]
另一方面 \(\operatorname{Hess}f=0\)，\(|\operatorname{Ric}|^2=n\lambda^2/a^2\)，故平方积分恰为相同数值。即使 \(\lambda<0\)，泛函也从负侧单调增加；其单调性没有预设标量曲率非负。
\end{solution}

\begin{exercise}\label{g3:ex:ricci-type-I-bound}
设闭 Ricci 流在有限时间 \(T\) 出现 I 型奇点，且 \(|\mathrm{Rm}|(t)\le C/(T-t)\)。对任意 \(t_j\uparrow T\)、\(Q_j>0\)，令 \(g_j(\sigma)=Q_jg(t_j+\sigma/Q_j)\)。证明在定义区间内
\[
|\mathrm{Rm}_{g_j(\sigma)}|\le
\frac{C}{Q_j(T-t_j)-\sigma}.
\]
若 \(Q_j=\max_M|\mathrm{Rm}|(t_j)\)，证明 \(Q_j(T-t_j)\) 同时有正的下界和有限上界。
\end{exercise}
\begin{solution}
缩放使曲率范数乘 \(Q_j^{-1}\)，因此直接代入原估计便得所列分式。I 型条件给 \(Q_j(T-t_j)\le C\)；\cref{g3:ex:ricci-curvature-rate}给 \(Q_j(T-t_j)\ge c_n>0\)。取子列可使这些数收敛到 \([c_n,C]\) 中的数，故归一化后的剩余终止时间有一致控制。这个实数子列的收敛并不等于度量序列的光滑收敛；后者还需要相应的局部几何估计与紧性。
\end{solution}

\begin{exercise}\label{g3:ex:ricci-normalized-nil}
对\cref{g3:prop:ricci-flow-nil-model}在紧商上的解，求保持初始体积的归一化系数、归一化时间以及三个度量分量。证明归一化后中心圆纤维长度仍趋于零，并求曲率范数随新时间的衰减阶。
\end{exercise}
\begin{solution}
记 \(L=1+3q_0t\)。体积为 \(V_0L^{1/6}\)，维数为三，故缩放因子为 \(a=L^{-1/9}\)。因此
\[
\tau=\int_0^t(1+3q_0s)^{-1/9}\,\mathrm ds
=\frac{3}{8q_0}(L^{8/9}-1),
\qquad
\widehat A=A_0L^{2/9},\quad
\widehat B=B_0L^{2/9},\quad
\widehat C=C_0L^{-4/9}.
\]
三系数乘积不变，验证体积归一化。中心圆纤维长度为原初长度乘 \(L^{-2/9}\)，即与 \((1+(8q_0/3)\tau)^{-1/4}\) 成正比，趋于零。未归一化曲率范数为一个正常数乘 \(q_0/L\)；缩放后再乘 \(a^{-1}=L^{1/9}\)，所以与 \((1+(8q_0/3)\tau)^{-1}\) 成正比。这给出固定总体积、曲率趋零而某些闭曲线仍缩短的具体演化。
\end{solution}

\begin{exercise}\label{g3:ex:ricci-cylinder-blowup}
在 \(S^2\times(\mathbb R/2\pi\mathbb Z)\) 上取 \(g_0=a_0g_{S^2(1)}+\ell^2\mathrm d\theta^2\)，其中 \(a_0,\ell>0\)。令 \(T=a_0/2\)，取 \(t_j\uparrow T\)、\(a_j=a_0-2t_j\)，并以 \(Q_j=a_j^{-1}\) 作抛物重缩放。用局部坐标证明重缩放流趋于 \(S^2\times\mathbb R\) 上的缩圆柱。
\end{exercise}
\begin{solution}
原流为 \(g(t)=(a_0-2t)g_{S^2(1)}+\ell^2\mathrm d\theta^2\)，终点处球面曲率发散。重缩放得到
\[
g_j(\sigma)=(1-2\sigma)g_{S^2(1)}
+\frac{\ell^2}{a_j}\mathrm d\theta^2,\qquad
-t_j/a_j\le\sigma<\frac12.
\]
固定圆上基点 \(\theta_0\)，在 \(|z|<R\) 上用 \(\theta=\theta_0+\sqrt{a_j}z/\ell\)。当 \(j\) 充分大，这个区间在圆上无自交，拉回度量恰为
\[
(1-2\sigma)g_{S^2(1)}+\mathrm dz^2.
\]
任意固定 \(R\) 和紧时间区间 \([\sigma_1,\sigma_2]\subset(-\infty,1/2)\) 最终均可使用同样的坐标。于是这些局部拉回度量连同全部导数都等于极限度量，直接得到所述光滑局部极限。圆的重缩放周长 \(2\pi\ell/\sqrt{a_j}\) 趋于无穷，解释了轴向的 \(\mathbb R\) 从何而来。
\end{solution}
